\documentclass{article}
\usepackage{readmasters}
\begin{document}

\origpage{189}
\section*{Ueber die Anwendung der Abelschen Functionen in der Geometrie.}

(Von Herrn \emph{A.~Clebsch} zu Giessen.)

Die Anwendungen, welche die elliptischen Functionen in der Geometrie bisher gefunden haben, sind, obwohl auf ganz wenige unter den möglichen Fällen beschränkt, dennoch schon zahlreich genug, um den allgemeinen Charakter derselben hervortreten zu lassen. Alle diese Anwendungen gehen von der Theilung aus, und zeigen, wie die Lösungen von gewissen in der Geometrie auftretenden algebraischen Gleichungen mit Hülfe der Theilung der elliptischen Functionen einfach dargestellt werden können. Dies ist um so mehr werth, als jene Gleichungen selbst ihrer verwickelten Form wegen im Allgemeinen sehr schwer aufzustellen sind; nicht minder aber deswegen, weil die eigenthümlichen Beziehungen und Gruppirungen der Wurzeln unter einander dabei aufs klarste hervortreten. Dass \emph{Steiner} in so vielen Fällen diesen Beziehungen nachzugehen gewusst hat, und dass viele von ihm gelöste Probleme durch die Theorie der elliptischen Functionen erst ihre wahre analytische Ausdrucksweise finden, ist nicht die geringste der merkwürdigen Thatsachen, welche uns bei dem Studium der Schriften dieses grossen Geometers mit stets wachsender Bewunderung erfüllen.

Aehnliche Anwendungen der Theorie der \emph{Abel}schen Functionen vorzuführen, ist der Zweck des gegenwärtigen Aufsatzes; und ich glaube, dass die darzulegenden wenigen Principien dazu dienen können, das Studium der algebraischen Curven in vieler Beziehung zu erleichtern. Dass solche Anwendungen bisher nicht versucht sind, obgleich wir seit sechs Jahren \emph{Riemann}s Theorie dieser Functionen besitzen, ist ohne Zweifel grossentheils den Schwierigkeiten zuzuschreiben, welche dem Verständnisse der betreffenden Abhandlungen noch immer entgegenstehen, und welche auch durch die neueren Bemühungen jüngerer Mathematiker nicht ganz gehoben sind. Ein anderer Grund liegt darin, dass man bei diesen Anwendungen zu derjenigen Anschauungsweise der Natur dieser Functionen zurückkehren muss, welche \emph{Jacobi} im $9^{\text{ten}}$ und $13^{\text{ten}}$ Bande dieses Journals aufgestellt hat, und welche sich bei Herrn \emph{Riemann} durch eine andere ersetzt findet, obwohl für den Uebergang

\origpage{190}
zu der \emph{Jacobi}schen Anschauungsweise das nöthige Material gegeben ist. Man kann sagen, dass bei \emph{Jacobi} die gegenseitige Abhängigkeit zweier Reihen von gleich viel Variabeln in den Vordergrund tritt, während es sich bei Herrn \emph{Riemann} vorzugsweise um eine Reihenentwicklung handelt. Ich werde im Folgenden immer von der Vorstellungsweise, wie \emph{Jacobi} sie benutzt, Gebrauch machen. Die von der Theorie der \emph{Abel}schen Functionen gemachten Ausführungen (von Herrn \emph{C.~Neumann}~*) und Herrn \emph{Prym}~**)), scheinen sogar zu zeigen, dass dieser \emph{Jacobi}sche Standpunkt den nothwendigen Dnrchgangspunkt bildet, und dass alle Entwicklungen zuerst in der von \emph{Jacobi} vorausgesagten Form auftreten, selbst wenn man, wie Herr \emph{Prym} thut, dieselbe zu vermeiden bestrebt ist.

Die geometrischen Anwendungen dieser Theorie, welche im Folgenden angebahnt werden, sind so zahlreich, dass ich nur einige Fälle näher auszuführen versucht habe. Diese Ausführungen beziehen sich namentlich auf die Theorie der ebenen Curven vierter Ordnung und enthalten unter anderen eine Lösung des Problems der Doppeltangenten, welche sich durch eine aus der Natur der Sache fliessende übersichtliche Gruppirung dieser merkwürdigen Geraden empfiehlt. Dabei ergeben sich fast von selbst eine Reihe von Sätzen, welche \emph{Hesse} und \emph{Steiner} im $49^{\text{ten}}$ Bande dieses Journals angegeben haben, und auf welche \emph{Hesse} seine elegante Deutung der zwischen den 28 Doppeltangenten eintretenden Beziehungen durch die Verbindungslinien von 8 Punkten im Raume gegründet hat. Es liegt in der Natur der hier angewandten Methoden, dass alle diese Sätze als ganz specielle Fälle von anderen, sehr allgemeinen, erscheinen, welche nicht bloss für ebene Curven beliebiger Ordnung, sondern auch für alle doppelt gekrümmten Curven in gleicher Weise ihr Geltung behalten.

\section*{§.~1.}
\subsection*{Zusammenhang einer Curve $n^{\text{ter}}$ Ordnung mit einer Classe von Abelschen Integralen.}

Jede algebraische Gleichung $F(s, z) = 0$, vermöge deren $s$ als Function von $z$ bestimmt ist, begründet nach Herrn \emph{Riemann} eine Classe von \emph{Abel}schen Integralen. Aus den überall endlich bleibenden Integralen, welche man dabei erhält, setzen sich die Argumente von $\theta$-Functionen auf lineare Weise zusammen. Nimmt man die $p$ Argumente einer $\theta$-Function als bekannt an,

*) Die Umkehrung der \emph{Abel}schen Integrale, Halle 1863.

**) Theoria nova ff. ultraellipticarum, Diss. Berlin 1863.

\origpage{191}
so entsprechen denselben $p$ Werthe von $z$, welche die oberen Grenzen der constituirenden Integrale bilden; und dieselben, oder irgend welche algebraische Functionen derselben, bestimmen sich als Wurzeln einer Gleichung $p^{\text{ten}}$ Grades, deren Coefficienten sich aus $\theta$-Functionen rational zusammensetzen. Die Zahl $p$ wird mit Hülfe der Formel
\[ p = \frac{\mathrm{w}}{2} - (n-1)\ \text{*)} \]
gefunden; dabei ist $n$ der Grad, bis zu welchem $s$ in der gegebenen Gleichung ansteigt, $\mathrm{w}$ aber die Zahl der Werthsysteme, für welche $F'(s) = 0$, ohne dass $F'(z) = 0$. Dabei wird vorausgesetzt, dass für diejenigen Werthsysteme, wo $F'(s) = 0$ und $F'(z) = 0$, nicht $F''(ss).F''(zz) - F''(sz)^{2} = 0$ sei.

Herr \emph{Riemann} legt die Gleichung $F(s, z) = 0$ zu Grunde in der Form, dass, wenn man die Gleichung nach Potenzen von $s$ ordnet, jeder Coefficient von gleich hoher Ordnung für $z$ sei. Für die geometrische Anwendung ist es zweckmässig, eine andere Form als allgemeine Grundform zu betrachten, nämlich diejenige, welche durch Substitution von $-\dfrac{a_{1}x_{1} + a_{2}x_{2} + a_{3}x_{3}}{b_{1}x_{1} + b_{2}x_{2} + b_{3}x_{3}}$ für $z$, $-\dfrac{\alpha_{1}x_{1} + \alpha_{2}x_{2} + \alpha_{3}x_{3}}{\beta_{1}x_{1} + \beta_{2}x_{2} + \beta_{3}x_{3}}$ für $s$ in eine homogene Function $n^{\text{ter}}$ Ordnung der drei Veränderlichen $x_{1}$, $x_{2}$, $x_{3}$ übergeführt werden kann, wo die $a$, $b$, $\alpha$, $\beta$ beliebige Constante bezeichnen.

Sei die so entstehende Gleichung
\[ f(x_{1}, x_{2}, x_{3}) = 0 \tag{1.} \]
die Gleichung einer Curve $n^{\text{ter}}$ Ordnung. Die Reduction dieser Gleichung auf die Form
\[ F(s, z) = 0 \tag{2.} \]
mit Hülfe der Gleichungen
\[
\left\{
\begin{aligned}
(a_{1}x_{1} + a_{2}x_{2} + a_{3}x_{3}) + z(b_{1}x_{1} + b_{2}x_{2} + b_{3}x_{3}) &= 0, \\
(\alpha_{1}x_{1} + \alpha_{2}x_{2} + \alpha_{3}x_{3}) + s(\beta_{1}x_{1} + \beta_{2}x_{2} + \beta_{3}x_{3}) &= 0
\end{aligned}
\right.
\tag{3.}
\]
entspricht dann der Darstellung der Curve als Durchschnitt entsprechender Strahlen der Büschel (3.), in denen $s$, $z$ die veränderlichen Parameter bilden. Untersuchen wir, wie gross die Zahl $\mathrm{w}$ hienach wird.

Wenn $F's = 0$, ohne dass $F'z = 0$, so hat die Gleichung $z$ zwei gleiche Wurzeln; zwei aufeinanderfolgenden Strahlen des ersten Büschels entspricht

*) Band 54 dieses Journals p.~129.

\origpage{192}
\emph{einer} des zweiten. Mit anderen Worten, ein Strahl des zweiten Büschels wird Tangente der Curve. Der Fall $F''s = 0$, $F'''s = 0$ etc. ist leicht zu berücksichtigen; die Tangente berührt dann in höherer Ordnung. Aber dies kann nur für besondere Werthe der $a$, $b$, $\alpha$, $\beta$ eintreten.

Wenn $F'(s) = 0$ und $F'(z) = 0$, so ist auch
\[ \frac{\partial f}{\partial x_{1}} = 0, \qquad \frac{\partial f}{\partial x_{2}} = 0, \qquad \frac{\partial f}{\partial x_{3}} = 0; \]
die Curve hat dann einen Doppelpunkt. Dass die Bedingung
\[ F''(ss).F''(zz) - F''(sz)^{2} = 0 \]
ausgeschlossen ist, bedeutet, dass kein Doppelpunkt ein Rückkehrspunkt werden solle~*).

Ist also die Zahl der Doppelpunkte gleich $d$, so hat man nach einer bekannten Formel für die Zahl der Tangenten,
\[ \mathrm{w} = n.(n-1) - 2d. \]
Daher
\[ p = \frac{(n-1)(n-2)}{2} - d. \tag{4.} \]
Dieses ist also die Zahl, welche die Classe der zu einer Curve $n^{\text{ter}}$ Ordnung gehörigen \emph{Abel}schen Functionen angiebt.

So entsprechen die elliptischen Functionen $(p = 1)$ den allgemeinen Curven dritter Ordnung, den Curven vierter Ordnung mit zwei Doppelpunkten etc.; die erste Classe der \emph{Abel}schen Transcendenten den Curven vierter Ordnung mit einem Doppelpunkt, oder denen fünfter Ordnung mit vier Doppelpunkten etc.; die folgende Classe $(p = 3)$ den allgemeinen Curven vierter Ordnung etc.

Nach der Theorie der \emph{Abel}schen Functionen findet man $z$ mit Hülfe einer Gleichung $p^{\text{ten}}$ Grades durch $\theta$-Functionen ausgedrückt. Die erste Gleichung (3.) lehrt also, wenn die $p$ Argumente der $\theta$ als bekannt vorausgesetzt werden, eine lineare Relation zwischen den Coordinaten eines Punktes kennen; das zu $z$ gehörige $s$, dessen Ausdruck man aus dem von $z$ und aus seinen nach den Argumenten der $\theta$-Functionen genommenen Differentialquotienten ableitet, liefert eine zweite, und man kann endlich die Coordinaten einer Curve so durch $\theta$-Functionen ausdrücken, wobei, neben der allgemeinen Theorie der

*) Diese Beschränkung lässt sich, wie es scheint, aufheben, indem man festsetzt, dass für jeden Rückkehrpunkt zugleich die Zahl der Tangenten um eine vermehrt werden solle.

\origpage{193}
betreffenden \emph{Abel}schen Functionen, nur die Auflösung einer einzigen Gleichung $p^{\text{ten}}$ Grades gefordert wird.

Die überall endlichen Integrale, aus denen die Argumente der $\theta$-Functionen sich zusammensetzen, kann man nach \emph{Riemann} (a.\ a.\ O. pag.~131.) und mit Hülfe der Gleichungen
\[ \frac{x_{2}\,dx_{3} - x_{3}\,dx_{2}}{\dfrac{\partial f}{\partial x_{1}}} = \frac{x_{3}\,dx_{1} - x_{1}\,dx_{3}}{\dfrac{\partial f}{\partial x_{2}}} = \frac{x_{1}\,dx_{2} - x_{2}\,dx_{1}}{\dfrac{\partial f}{\partial x_{3}}} \]
in die Form bringen:
\[ \int \Theta.\frac{\Sigma\pm c_{1}x_{2}\,dx_{3}}{c_{1}\dfrac{\partial f}{\partial x_{1}} + c_{2}\dfrac{\partial f}{\partial x_{2}} + c_{3}\dfrac{\partial f}{\partial x_{3}}}, \]
wobei $\Theta$ eine ganze homogene Function der Ordnung $n-3$ ist, welche für die Doppelpunkte der Curve verschwindet, und wo zwischen den $x$ die Gleichung $f = 0$ bestehend angesehen wird. Die $c$, welche der Symmetrie wegen eingeführt sind (vergl. \emph{Aronhold}, Monatsber. der Berl. Acad., Sitzung vom 25. April 1861), haben auf das Integral gar keinen Einfluss, indem die Grösse unter dem Integralzeichen in Wahrheit von denselben unabhängig ist. Die zweckmässige Wahl der unteren Grenzen so wie der in den $\Theta$ auftretenden Coefficienten hat Herr \emph{Riemann} gelehrt; ich werde beides so voraussetzen, wie es in der zweiten Abtheilung der angeführten Abhandlung gegeben ist.

\section*{§.~2.}
\subsection*{Anwendung des Abelschen Satzes auf den Durchschnitt zweier Curven.}

Der \emph{Abel}sche Satz (Band 4 pag.~200 dieses Journals) kann nun, mit Zugrundelegung homogener Gleichungen statt der von \emph{Abel} benutzten, in einer Weise ausgesprochen werden, welche wichtige Folgerungen gestattet. Ich werde den Beweis hinzufügen, wie derselbe sich bei der angegebenen Modification gestaltet.

Es sei $f = 0$ die Gleichung einer Curve $n^{\text{ter}}$ Ordnung, $\varphi = 0$ die einer anderen Curve $m^{\text{ter}}$ Ordnung. Es sei ferner $\Theta$ eine beliebige rationale homogene Function $(n-3)^{\text{ter}}$ Ordnung. Betrachten wir die Summe
\[ \Sigma\int \Theta.\frac{\Sigma\pm c_{1}x_{2}\,dx_{3}}{c_{1}\dfrac{\partial f}{\partial x_{1}} + c_{2}\dfrac{\partial f}{\partial x_{2}} + c_{3}\dfrac{\partial f}{\partial x_{3}}}, \]

\origpage{194}
ausgedehnt über alle diejenigen Punkte als obere Grenzen der Integrale, welche das Schnittsystem von $f = 0$, $\varphi = 0$ bilden.

Unter dem Integralzeichen genügen die $x$ keineswegs der Gleichung $\varphi = 0$, sofern dieselbe eine bestimmte Curve repräsentirt; aber wir können die Gleichung $\varphi = 0$ bestehend annehmen, indem wir ihre Coefficienten als Variable einführen, welchen erst in den Grenzen die betreffenden gegebenen Werthe beizulegen sind.

Multiplicirt man nun unter dem Integralzeichen Zähler und Nenner mit $\Sigma\pm k_{1}\dfrac{\partial f}{\partial x_{2}}\dfrac{\partial \varphi}{\partial x_{3}}$, wo die $k$ beliebige Grössen sind, so findet man im Zähler die Determinante:
\[
\begin{vmatrix}
\Sigma k_{i}c_{i}, & \Sigma k_{i}x_{i}, & \Sigma k_{i}\,dx_{i} \\[4pt]
\Sigma c_{i}\dfrac{\partial f}{\partial x_{i}}, & \Sigma x_{i}\dfrac{\partial f}{\partial x_{i}}, & \Sigma \dfrac{\partial f}{\partial x_{i}}\,dx_{i} \\[4pt]
\Sigma c_{i}\dfrac{\partial \varphi}{\partial x_{i}}, & \Sigma x_{i}\dfrac{\partial \varphi}{\partial x_{i}}, & \Sigma \dfrac{\partial \varphi}{\partial x_{i}}\,dx_{i}
\end{vmatrix}.
\]
Von den Elementen dieser Determinante verschwinden drei; das letzte wird gleich $-\delta\varphi$, wenn $\delta\varphi$ dasjenige bedeutet, was aus $\varphi$ durch Differentiation der Coefficienten entsteht; und man hat endlich als Werth der Determinante:
\[ \Sigma c_{i}\frac{\partial f}{\partial x_{i}}.\Sigma k_{i}x_{i}.\delta\varphi, \]
so dass die gesuchte Summe übergeht in
\[ \Sigma\int \Theta.\frac{\Sigma k_{i}x_{i}}{\Sigma\pm k_{1}\dfrac{\partial f}{\partial x_{2}}\dfrac{\partial \varphi}{\partial x_{3}}}.\delta\varphi. \]
Hier sind jetzt die Incremente der einzelnen Coefficienten von $\varphi$ in rationale Functionen der Coordinaten der verschiedenen Punkte multiplicirt, welche den Gleichungen $f = 0$, $\varphi = 0$ zusammen genügen. Sämmtliche Punkte kommen in diesen Functionen symmetrisch vor; man kann dieselben daher als rationale Functionen der Coefficienten von $f$ und $\varphi$ darstellen; und indem man nach letzteren integrirt, erhält man, neben rationalen Functionen der Coefficienten nur logarithmische und Kreisfunctionen von solchen. Dies ist der \emph{Abel}sche Satz.

Bemerken wir nun insbesondere, dass, wenn $\Theta$ zugleich eine \emph{ganze} Function ist, die Summe der unter dem Integralzeichen stehenden Functionen immer verschwindet. Indem man nämlich das von \emph{Jacobi} (Band 14 dieses

\origpage{195}
Journals p.~286) ausgesprochene Theorem auf homogene Gleichungen anwendet, und zugleich bemerkt, dass, wenn $f = 0$, $\varphi = 0$, auch:
\[ \frac{\dfrac{\partial f}{\partial x_{2}}\dfrac{\partial \varphi}{\partial x_{3}} - \dfrac{\partial f}{\partial x_{3}}\dfrac{\partial \varphi}{\partial x_{2}}}{x_{1}} = \frac{\dfrac{\partial f}{\partial x_{3}}\dfrac{\partial \varphi}{\partial x_{1}} - \dfrac{\partial f}{\partial x_{1}}\dfrac{\partial \varphi}{\partial x_{3}}}{x_{2}} = \frac{\dfrac{\partial f}{\partial x_{1}}\dfrac{\partial \varphi}{\partial x_{2}} - \dfrac{\partial f}{\partial x_{2}}\dfrac{\partial \varphi}{\partial x_{1}}}{x_{3}}, \]
so erhält man die Formel:
\[ \Sigma\frac{\Omega(k_{1}x_{1} + k_{2}x_{2} + k_{3}x_{3})}{\Sigma\pm k_{1}\dfrac{\partial f}{\partial x_{2}}\dfrac{\partial \varphi}{\partial x_{3}}} = 0, \]
wo $\Omega$ eine ganze homogene Function der Ordnung $m+n-3$ ist, und wo die Summe sich auf alle Werthsysteme bezieht, welche den Gleichungen $f = 0$, $\varphi = 0$ gleichzeitig genügen. Da nun $\Theta.\delta\varphi$ von der Form $\Omega$ ist, sobald $\Theta$, welches von der $(n-3)^{\text{ten}}$ Ordnung war, eine \emph{ganze} Function bedeutet, so wird die obige Integralsumme dann constant, d.\ h.\ von den Coefficienten von $\varphi$ unabhängig

Man hat also folgenden Satz:

\emph{Bedeutet $f = 0$ die Gleichung einer Curve $n^{\text{ter}}$ Ordnung, und wird das Integral}
\[ \int \Theta.\frac{\Sigma\pm c_{1}x_{2}\,dx_{3}}{c_{1}\dfrac{\partial f}{\partial x_{1}} + c_{2}\dfrac{\partial f}{\partial x_{2}} + c_{3}\dfrac{\partial f}{\partial x_{3}}}, \]
\emph{in welchem $\Theta$ eine homogene Function $(n-3)^{\text{ter}}$ Ordnung ist, so gebildet, dass man die $x$ durch die Gleichung $f = 0$ von einander abhängig macht, so ist die Summe aller Werthe des Integrals für die Schnittpunkte von $f = 0$ mit einer anderen Curve $\varphi = 0$ ($m^{\text{ter}}$ Ordnung) als Aggregat von rationalen und von Logarithmen rationaler Functionen der Coefficienten von $\varphi$ darstellbar; insbesondere aber ist diese Summe von den Coefficienten von $\varphi$ unabhängig, sobald $\Theta$ zugleich eine ganze Function ist.}

Da $\Theta$, wenn es eine ganze Function ist, noch $\dfrac{n-1.n-2}{2}$ willkürliche Coefficienten enthält, so erhält man hieraus ebensoviel einzelne Gleichungen, welche Summen von $mn$ Integralen Constanten gleich geben. Diese $\dfrac{n-1.n-2}{2}$ Gleichungen enthalten in keiner Weise die Coefficienten von $\varphi$; sie sind daher überhaupt die Bedingungen dafür, dass $mn$ Punkte einer Curve $n^{\text{ter}}$ Ordnung auf einer Curve $m^{\text{ter}}$ Ordnung liegen. Ist $m \geqq n$, so ist bekanntlich die An-

\origpage{196}
zahl der hierzu erforderlichen Bedingungen genau gleich $\dfrac{n-1.n-2}{2}$~*); ist aber $m$ kleiner als $n$, so ist die Anzahl der erforderlichen Bedingungen:
\[ mn - \frac{m.m+3}{2}, \]
und da
\[ \left(mn - \frac{m.m+3}{2}\right) - \frac{n-1.n-2}{2} = -\frac{n-m-1.n-m-2}{2}, \]
so ist auch noch für $m = n-1$, $m = n-2$ die Anzahl der erhaltenen Gleichungen der Anzahl der erforderlichen Bedingungen gleich. Ist $m < n-2$, so wird die Zahl der Bedingungsgleichungen grösser; dieselben können aber nicht aufhören mit einander verträglich zu sein, und die hinreichende Anzahl von Bestimmungen zu ersetzen.

Wenn die betrachtete Curve keinen Doppelpunkt besitzt, also im Allgemeinen, bestehen die durch das Theorem gegebenen Gleichungen nur zwischen überall endlich bleibenden Integralen, wie sie bei der Bildung der Argumente der $\Theta$-Functionen gebraucht werden. Dieser allgemeine Fall soll zunächst behandelt werden. Existiren Doppelpunkte, so sind in ebensoviel Gleichungen die Integrale nicht mehr unter diese Categorie gehörig.

\section*{§.~3.}
\subsection*{Bestimmung der Constanten des Abelschen Satzes, und Umkehrung des Satzes über die Schnittpunkte algebraischer Curven.}

Setzen wir im Folgenden voraus, die Curve besitze keinen Doppelpunkt. Aus diesem Fall können alle anderen als specielle abgeleitet werden. Wir haben $p = \dfrac{n-1.n-2}{2}$; die überall endlich bleibenden Integrale bezeichnen wir durch
\[ u_{1}, \quad u_{2}, \quad \ldots \quad u_{p}, \]
und unterscheiden das verschiedenen Punkten angehörige durch obere Indices. Die $u$ sollen so bestimmt sein, wie sie Herr \emph{Riemann} zur Bildung der Argumente der $\theta$-Functionen verwendet; dadurch sind einerseits die unteren Grenzen bestimmt, andererseits ist festgesetzt, welche Veränderungen die $u$ durch Aenderung des Integrationsweges erfahren, dass sie nämlich bei gleichbleibender oberer Grenze dadurch nur übergeführt werden können in

*) \emph{Jacobi}, dieses Journal Band 15. pag.~285.

\origpage{197}
das System:
\[ u_{1} + A_{1}, \quad u_{2} + A_{2}, \quad \ldots \quad u_{p} + A_{p}, \]
wo, wie immer im Folgenden, die $A$ ein System der Form
\[
\left\{
\begin{aligned}
A_{1} &= m_{1}i\pi + a_{11}q_{1} + a_{12}q_{2} + \cdots + a_{1p}q_{p}, \\
A_{2} &= m_{2}i\pi + a_{21}q_{1} + a_{22}q_{2} + \cdots + a_{2p}q_{p}, \\
&\ \cdot\quad\cdot\quad\cdot\quad\cdot\quad\cdot\quad\cdot\quad\cdot\quad\cdot \\
A_{p} &= m_{p}i\pi + a_{p1}q_{1} + a_{p2}q_{2} + \cdots + a_{pp}q_{p}
\end{aligned}
\right.
\tag{5.}
\]
bezeichnen. In diesen Ausdrücken sind die $a$ durch die Curvengleichung gegeben, die $m$, $q$ aber bedeuten ganze Zahlen.

Die Gleichungen des Theorems gehen dann in folgende Form über:
\[
\left\{
\begin{aligned}
u'_{1} + u_{1}^{(2)} + \cdots + u_{1}^{(mn)} &\equiv \gamma_{1}, \\
u'_{2} + u_{2}^{(2)} + \cdots + u_{2}^{(mn)} &\equiv \gamma_{2}, \\
&\ \cdot\quad\cdot\quad\cdot\quad\cdot\quad\cdot\quad\cdot \\
u'_{p} + u_{p}^{(2)} + \cdots + u_{p}^{(mn)} &\equiv \gamma_{p}.
\end{aligned}
\right.
\tag{6.}
\]
Da die Curve $m^{\text{ter}}$ Ordnung aus $m$ Geraden bestehen kann, so muss man jede dieser Gleichungen dadurch bilden können, dass man die $m$ entsprechenden Gleichungen addirt, welche für $m = 1$ bestehen. Man hat daher, wenn $c_{1}$, $c_{2}$, $\ldots$ $c_{p}$ die zu $m = 1$ gehörigen Constanten $\gamma$ bedeuten, sogleich
\[ \gamma_{1} = mc_{1}, \quad \gamma_{2} = mc_{2}, \quad \ldots \quad \gamma_{p} = mc_{p}. \]
Die Constanten $c$ sind dann nicht blos von den Coefficienten sondern auch von dem Grade der Curve $\varphi = 0$ unabhängig. Ich werde jetzt zeigen, dass die $c$ sämmtlich gleich Null gesetzt werden können.

Hiezu führt die Betrachtung des Falles $m = n-3$. Herr \emph{Riemann} nämlich hat a.\ a.\ O. p.~148 bewiesen, dass ein System der obigen Art congruent mit lauter Nullen ist, sobald es ausgedehnt wird über die gemeinschaftlichen Lösungen von $f = 0$ mit denjenigen Gleichungen, deren linker Theil aus den Zählern der in den $u$ integrirten Functionen linear zusammengesetzt ist, etwaige Doppelpunkte ausgeschlossen, welche hier von vorn herein nicht vorkommen. Nun ist eine lineare Function der Zähler der integrirten Ausdrücke eine allgemeine homogene Function $(n-3)^{\text{ter}}$ Ordnung $(\Theta)$; für $m = n-3$ ist also $(n-3)c_{h} = 0$, also immer $c_{h} = 0$. Man hat daher folgenden Satz:

\emph{Bildet man die überall endlichen Integrale mit den von Herrn Riemann angegebenen unteren Grenzen, und bildet die Summen solcher In-}

\origpage{198}
\emph{tegrale für die Durchschnittspunkte der Curve $f = 0$ mit einer anderen, so ist das System der erhaltenen Integralsummen jederzeit einem System von Nullen congruent.}

Dieser Satz gestattet die Umkehrung und lautet dann so:

\emph{Sobald die Summen der über $m.n$ Punkte ausgedehnten überall endlichen Integrale Null sind, liegen die $m.n$ Punkte im Durchschnitt einer Curve $m^{\text{ter}}$ Ordnung mit der Curve $n^{\text{ter}}$ Ordnung.}

Die Differentialgleichungen nämlich:
\[
\begin{aligned}
du'_{1} + du_{1}^{(2)} + \cdots + du_{1}^{(mn)} &= 0, \\
du_{2}^{(1)} + du_{2}^{(2)} + \cdots + du_{2}^{(mn)} &= 0, \\
&\ \cdot\quad\cdot\quad\cdot\quad\cdot\quad\cdot\quad\cdot \\
du_{p}^{(1)} + du_{p}^{(2)} + \cdots + du_{p}^{(mn)} &= 0
\end{aligned}
\]
werden durch $p$ algebraische Gleichungen, welche der \emph{Abel}sche Satz giebt, vollständig integrirt. Die $p$ Integrale enthalten $p$ willkürliche Constanten. Da die Differentialgleichungen ferner durch solche algebraische Gleichungen integrirt werden, welche aussagen, dass die $mn$ Punkte auf einer Curve $m^{\text{ter}}$ Ordnung liegen, so müssen diese letzteren aus jenen durch eine specielle Wahl der willkürlichen Constanten ableitbar sein, und wenn man solche Werthe der Constanten festhält, so müssen die algebraischen Integrale eben jene Gestalt haben, in welcher sie die $mn$ Punkte als den Durchschnittspunkt einer Curve $m^{\text{ter}}$ Ordnung mit der gegebenen charakterisiren. Die erforderliche Constantenbestimmung ist aber oben geleistet worden.

\section*{§.~4.}
\subsection*{Berührungscurven, welche durch eine Anzahl auf der Curve gegebener Punkte bestimmt sind.}

Mit Hülfe der entwickelten Sätze ergiebt sich leicht die Lösung folgender Aufgabe:

\emph{Es sei $m \geqq n-2$; auf der Curve $n^{\text{ter}}$ Ordnung sind $mn-pr$ Punkte beliebig gegeben; man soll durch sie eine Curve $m^{\text{ter}}$ Ordnung legen, welche die Curve $n^{\text{ter}}$ Ordnung in $p$ Punkten $r$-punktig berührt.}

Diese Aufgabe ist lösbar, wenn $m \geqq n-2$. Sie ist auch bestimmt, wenn $m = n-2$ oder $m = n-1$. Denn in beiden Fällen ist die Anzahl der gestellten Bedingungen, nämlich $mn - pr + p(r-1)$ gleich der Zahl der in der Curvengleichung enthaltenen Constanten. Ist $m \geqq n$, so ist die Zahl der Bedingungsgleichungen $(mn - p)$ nach einem bereits erwähnten Satze eben aus-

\origpage{199}
reichend, um das Punktsystem zu bestimmen, welches die Curve $m^{\text{ter}}$ Ordnung mit der Curve $n^{\text{ter}}$ Ordnung gemein haben muss; aber es erfüllt sodann die gestellten Bedingungen jede Curve $m^{\text{ter}}$ Ordnung, welche durch die gefundenen Punkte hindurchgeht.

Setzen wir der Kürze wegen $mn - pr = \lambda$. Von den Schnittpunkten beider Curven sind dann $\lambda$ gegeben, von den übrigen $pr$ sollen je $r$ zusammenfallen. Die Summen entsprechender Integrale, welche den gesuchten Berührungspunkten zugehören, seien $v_{1}$, $v_{2}$, $\ldots$ $v_{p}$; die Gleichungen (6.) nehmen dann die Gestalt an:
\[
\begin{aligned}
rv_{1} &\equiv -(u'_{1} + u_{1}^{(2)} + \cdots + u_{1}^{(\lambda)}), \\
rv_{2} &\equiv -(u'_{2} + u_{2}^{(2)} + \cdots + u_{2}^{(\lambda)}), \\
&\ \cdot\quad\cdot\quad\cdot\quad\cdot\quad\cdot\quad\cdot \\
rv_{p} &\equiv -(u'_{p} + u_{p}^{(2)} + \cdots + u_{p}^{(\lambda)}).
\end{aligned}
\]
Die rechten Theile dieser Congruenzen sind gegeben; durch Division mit $r$ findet man, indem man der Allgemeinheit wegen rechts ein willkürliches System der $A$, durch $r$ dividirt, hinzufügt:
\[
\left\{
\begin{aligned}
v_{1} &\equiv -\frac{u'_{1} + u_{1}^{(2)} + \cdots + u_{1}^{(\lambda)}}{r} + \frac{A_{1}}{r}, \\
v_{2} &\equiv -\frac{u'_{2} + u_{2}^{(2)} + \cdots + u_{2}^{(\lambda)}}{r} + \frac{A_{2}}{r}, \\
&\ \cdot\quad\cdot\quad\cdot\quad\cdot\quad\cdot\quad\cdot \\
v_{p} &\equiv -\frac{u'_{p} + u_{p}^{(2)} + \cdots + u_{p}^{(\lambda)}}{r} + \frac{A_{p}}{r}.
\end{aligned}
\right.
\tag{7.}
\]
Bildet man jetzt die $\theta$, deren Argumente eben diese $v$ sind, und die zugehörige Gleichung $p^{\text{ten}}$ Grades, so findet man durch Auflösung die $p$ Punkte, welche in den $v$ als obere Grenzen vorkommen, d.\ h.\ die $p$ Berührungspunkte, wodurch die Aufgabe gelöst ist.

Die $A$ enthalten die $2p$ Zahlen $m$, $q$, denen noch alle Werthe von 0 bis $r-1$ beigelegt werden können. \emph{Die Anzahl aller Lösungen des Problems ist daher $r^{2p}$.}

Die $r^{2p}$ Systeme von je $p$ Berührungspunkten haben eine sehr merkwürdige Lage gegen einander. Unterscheidet man $r$ solcher Systeme, von denen auch mehrere gruppenweise zusammenfallen können, durch beigesetzte obere Indices, und addirt die entsprechenden Gleichungen (7.), so erhält man

\origpage{200}
für $k = 1, 2, \ldots p$:
\[ v_{k}^{(1)} + v_{k}^{(2)} + \cdots + v_{k}^{(r)} + u_{k}^{(1)} + u_{k}^{(2)} + \cdots + u_{k}^{(\lambda)} \equiv \frac{A_{k}^{(1)} + A_{k}^{(2)} + \cdots + A_{k}^{(r)}}{r}. \tag{8.} \]
Bestimmt man also die in den $A$ vorkommenden Zahlen $m$, $q$ so, dass
\[
\left\{
\begin{aligned}
m_{k}^{(1)} + m_{k}^{(2)} + \cdots + m_{k}^{(r)} &\equiv 0 \quad (\text{mod. } r), \\
q_{k}^{(1)} + q_{k}^{(2)} + \cdots + q_{k}^{(r)} &\equiv 0 \quad (\text{mod. } r),
\end{aligned}
\right.
\tag{9.}
\]
so kann man statt der rechten Seite von (8.) Null setzen, und es findet sich also, dass die $r$ Systeme von Berührungspunkten mit den gegebenen Punkten in einer Curve $m^{\text{ter}}$ Ordnung liegen. Aus (9.) sieht man, dass von den $r$ Systemen $r-1$ beliebig gewählt werden können, wodurch dann das $r^{\text{te}}$ sich bestimmt; und man hat also den Satz:

\emph{Legt man eine Curve $m^{\text{ter}}$ Ordnung durch die gegebenen Punkte und durch $r-1$ Systeme von Berührungspunkten, so geht dieselbe noch durch ein $r^{\text{tes}}$ System. Von ersteren können mehrere gruppenweise zusammenfallen; es kann aber auch geschehen, dass letzteres mit einem der ersteren zusammenfällt; was insbesondere bei $r = 2$ immer eintritt.}

Es braucht kaum erwähnt zu werden, dass, wenn $\mu$ Systeme zusammenfallen, in den Punkten des Systems die neue Curve mit der gegebenen eine $\mu$-punktige Berührung hat.

\section*{§.~5.}
\subsection*{Systeme von Berührungscurven, welche durch feste Punkte der Curve gehen.}

Lassen wir in dem vorigen Problem von den gegebenen Punkten, deren Anzahl grösser als $\mu.r$ vorausgesetzt werden mag, $\mu$ mal $r$ Punkte zusammenfallen. Nehmen wir sodann nur die übrigen $mn - (p+\mu)r$ Punkte als gegeben an. Wir haben es dann mit der unbestimmten Aufgabe zu thun:

\emph{Eine Curve $m^{\text{ter}}$ Ordnung soll durch $mn - (p+\mu)r$ auf einer Curve $n^{\text{ter}}$ Ordnung gegebene Punkte gehen, und diese Curve ausserdem in $p+\mu$ Punkten $r$-punktig berühren.}

Diese Aufgabe ist, wenn $m \geqq n-2$, wie die vorige lösbar, und zwar sind $\mu$ Berührungspunkte noch willkürlich, die $p$ übrigen aber dann bestimmt. Durch die gegebenen $mn - (p+\mu)r$ Punkte lassen sich also noch unendlich viele Curven der gesuchten Art legen; wenn $\mu$ Berührungspunkte willkürlich gewählt sind, noch immer $r^{2p}$.

\origpage{201}
Herr \emph{Hesse} hat den Begriff eines \emph{Systems} von Berührungscurven eingeführt (Band 49 dieses Journals) und für die Theorie der Curven dritter und vierter Ordnung davon mannigfachen Gebrauch gemacht. Geht man nämlich von einer bestimmten Berührungscurve aus, und lässt dieselbe sich stetig verändern, so durchlaufen die Berührungspunkte allmälig alle nur möglichen Combinationen. Aber es ist keineswegs möglich, durch solche stetige Veränderung zu \emph{allen} Berührungscurven zu gelangen; sondern man muss zu diesem Zwecke diejenigen Berührungscurven, welche eine möglichst grosse Anzahl von beliebig gewählten Berührungspunkten gemeinsam haben, gleichzeitig variiren lassen. Aus jeder dieser Curven entsteht dann ein System von Berührungscurven, deren keines mit einem anderen eine Curve gemein hat.

Man bemerkt nun, dass die verschiedenen Systeme von Berührungscurven bei der vorliegenden Darstellung durch ein sehr einfaches Kennzeichen charakterisirt sind. Jedem System entspricht nämlich ein bestimmtes System der $A$; denn diese Systeme sind völlig discret, und es kann daher niemals geschehen, dass man von einer Curve, welche einem System der $A$ entspricht, zu Curven welche einem anderen System entsprechen, einen stetigen Uebergang durch lauter Berührungscurven hindurch herstelle. Man darf daher folgenden Satz aussprechen:

\emph{Man kann unendlich viele Curven $m^{\text{ter}}$ Ordnung, wenn $m \geqq n-2$, bestimmen, welche eine Curve $n^{\text{ter}}$ Ordnung in $\dfrac{n-1.n-2}{2} + \mu$ Punkten $r$-punktig berühren, während der Rest der Schnittpunkte gegeben ist. Diese Curven theilen sich in $r^{2p}$ verschiedene Systeme ein, dergestalt dass es nicht möglich ist, von einer Curve eines Systems durch lauter Berührungscurven hindurch zu einer Curve eines anderen Systems stetig überzugehen.}

Bezeichnen wir nun die den $p+\mu$ Berührungspunkten entsprechenden Integrale durch
\[ w_{k}^{(1)}, \quad w_{k}^{(2)}, \quad \ldots \quad w_{k}^{(p+\mu)}, \]
die Summen der den festen Punkten entsprechenden Integrale durch $U_{k}$, so haben wir:
\[ w_{k}^{(1)} + w_{k}^{(2)} + \cdots + w_{k}^{(p+\mu)} \equiv \frac{-U_{k} + A_{k}}{r}. \]
Betrachten wir jetzt $r$ Curven desselben Systems, und addiren die entsprechenden Gleichungen, was links durch das Zeichen $\mathrm{S}$ angedeutet werden mag,

\origpage{202}
so kommt:
\[ \mathrm{S}(w_{k}^{(1)} + w_{k}^{(2)} + \cdots + w_{k}^{(p+\mu)}) \equiv -U_{k}, \]
d.\ h.

\emph{Die Berührungspunkte von irgend $r$ Berührungscurven desselben Systems liegen jederzeit auf einer Curve $m^{\text{ter}}$ Ordnung, welche durch die gegebenen festen Punkte hindurchgeht.}

Ausserdem hat man, wie im vorigen §., den Satz:

\emph{Legt man eine Curve $m^{\text{ter}}$ Ordnung durch die gegebenen Punkte und durch $r-1$ Systeme von je $p+\mu$ Berührungspunkten, welche gleichen oder verschiedenen Systemen von Berührungscurven angehören, so geht dieselbe immer noch durch ein $r^{\text{tes}}$ System.}

Der vorige Satz kann als specieller Fall dieses letzteren aufgefasst werden, und für $r = 2$, bei überall zweipunktiger Berührung, geht dieser geradezu in jenen über.

\section*{§.~6.}
\subsection*{Systeme von Curven, welche die gegebene Curve in einer gewissen Ordnung berühren, wo sie derselben begegnen.}

Wenn $mn$ durch $r$ theilbar ist, so treten insbesondere Berührungscurven auf, wie diejenigen, welche Herr \emph{Hesse} a.\ a.\ O. untersucht hat, welche die Curve $n^{\text{ter}}$ Ordnung überall $r$-punktig berühren, wo sie derselben begegnen. Es ist dann $mn = (p+\mu)r$; die entstehenden Systeme sind allein von der Curve $n^{\text{ter}}$ Ordnung selbst abhängig, nicht mehr, wie im vorigen, ausserdem durch besondere auf ihr beliebig gewählte Punkte charakterisirt. Die übrigen im vorigen §. angeführten Sätze fahren hier fort zu gelten; aber die Anzahl der Systeme kann unter Umständen eine andere werden als die oben angegebene.

Nehmen wir nämlich an, es habe $m$ mit $r$ einen gemeinschaftlichen Factor, so dass
\[ m = m's, \qquad r = r's, \]
wo $m'$, $r'$ relative Primzahlen seien. So oft nun alle in den $A$ vorkommende Zahlen $m$, $q$ den Factor $s$ enthalten, hat man
\[ \frac{A_{k}}{r} = \frac{A'_{k}}{r'}, \]
wo die $A'$ Ausdrücke nach Art der $A$ sind. Dann ist immer, wenn für die Berührungspunkte die betreffenden Integrale durch $u_{k}^{(1)}$, $u_{k}^{(2)}$, $\ldots$ $u_{k}^{\left(\frac{mn}{r}\right)} = u_{k}^{\left(\frac{m'n}{r'}\right)}$

\origpage{203}
bezeichnet werden:
\[ u_{k}^{(1)} + u_{k}^{(2)} + \cdots + u_{k}^{\left(\frac{m'n}{r'}\right)} \equiv \frac{A'_{k}}{r'}, \]
d.\ h.\ das System von Berührungscurven besteht in Wahrheit aus Curven der Ordnung $m'$, welche die gegebene Curve überall $r'$-punktig berühren, und welche, $s$ mal gezählt, als Systeme $m^{\text{ter}}$ Ordnung mit $r$-punktiger Berührung angesehen werden können. Indem man von diesen abstrahirt, und bemerkt, dass diese Ausnahmsysteme der $A$ an der Zahl $r'^{2p}$ sind, gelangt man zu dem Satz:

\emph{Wenn $mn = (p+\mu)r$, $m \geqq n-2$, so giebt es Systeme von Curven $m^{\text{ter}}$ Ordnung, welche die Curve $n^{\text{ter}}$ Ordnung in $p+\mu$ Punkten $r$-punktig berühren, von denen $\mu$ beliebig sind. Ist $m = m's$, $r = r's$, wo $m'$, $r'$ relative Primzahlen, so ist die Anzahl der Systeme $r^{2p} - r'^{2p}$; nur wenn $m$ und $r$ relative Primzahlen sind, ist die Anzahl derselben gleich $r^{2p}$.}

Hieran knüpft sich eine Anzahl weiterer Sätze über diejenigen Fälle, in denen die Berührungspunkte mehrerer Berührungscurven zusammen den vollständigen Durchschnitt der Curve $n^{\text{ter}}$ Ordnung mit einer anderen Curve bilden. Solche Sätze existiren auch noch, wenn man nicht Berührungscurven gleicher Ordnung betrachtet. Es ist nämlich, wenn $m'$, $r'$; $m''$, $r''$; $\ldots$ $m^{(s)}$, $r^{(s)}$ die den verschiedenen zusammen betrachteten Berührungscurven entsprechenden Zahlen bezeichnen, und wenn die zugehörigen Systeme der $A$ in ähnlicher Weise bezeichnet werden, nur nöthig, dass
\[ \frac{m'}{r'} + \frac{m''}{r''} + \cdots + \frac{m^{(\mu)}}{r^{(\mu)}} = h. \]
eine ganze Zahl ist, und dass die Gleichungen stattfinden:
\[ \frac{A'_{k}}{r'} + \frac{A''_{k}}{r''} + \cdots + \frac{A_{k}^{(\mu)}}{r^{(\mu)}} \equiv 0, \]
damit die Berührungspunkte der Curven von den Ordnungen $m'$, $m''$, $\ldots$ $m^{(\mu)}$ auf einer Curve der $h^{\text{ten}}$ Ordnung liegen.

Ich begnüge mich diese Quelle von Sätzen hier angegeben zu haben.

\section*{§.~7.}
\subsection*{Berührung einer Curve vierter Ordnung durch Curven dritter Ordnung, vierpunktig in drei Punkten; und Berührung einer Curve sechster Ordnung durch Curven fünfter Ordnung, dreipunktig in 10 Punkten.}

Die im Vorigen behandelten Probleme sind wesentlich unbestimmt. Es giebt aber Classen von völlig bestimmten Berührungscurven. Für $m \geqq n$ treten

\origpage{204}
dieselben immer ein, wenn $mn = pr$, oder $\dfrac{m}{r} = \dfrac{p}{n}$. Hat nun $n$ nicht die Form $4h+2$, so haben $p$ und $n$ keinen gemeinschaftlichen Factor; man hat daher als allgemeinste Annahme $m = p.s$, $r = n.s$. Nur wenn $n$ die Form $4h+2$ hat, haben $p$ und $s$ den gemeinsamen Factor 2. In diesem Falle muss man daher setzen $m = \dfrac{p}{2}.s$, $r = \dfrac{n}{2}.s$. Und verbindet man dies mit den Bestimmungen des vorigen §., so ergiebt sich folgender Satz:

\emph{Ist $s$ eine beliebige Zahl, und $n$ nicht von der Form $4h+2$, so giebt es immer $n^{2p}(s^{2p}-1)$ Curven der Ordnung $p.s$, welche die Curve $n^{\text{ter}}$ Ordnung in $p$ Punkten $n.s$-punktig berühren. Ist dagegen $n$ von der Form $4h+2$, so giebt es $(\tfrac{1}{2}n)^{2p}.(s^{2p}-1)$ Curven der Ordnung $\tfrac{1}{2}p.s$, welche die Curve $n^{\text{ter}}$ Ordnung in $p$ Punkten $\tfrac{1}{2}n.s$-punktig berühren.}

Dieser Satz muss scheinbar zwei Ausnahmen erfahren, bei $s = 1$ und $n = 4$ oder $n = 6$, weil in diesen Fällen $m < n$ wird. Aber in beiden Fällen ist der Satz dennoch richtig, wovon man sich in folgender Weise überzeugt. Die Anzahl von Bedingungen, welchen eine in $\dfrac{mn}{r}$ Punkten $r$-punktig berührende Curve unterworfen ist, findet man gleich $\dfrac{mn}{r}(r-1)$. Damit hiedurch die Curve $m^{\text{ter}}$ Ordnung $(m < n)$ vollständig bestimmt sei, muss man die Gleichung haben:
\[ \frac{mn(r-1)}{r} = \frac{m.m+3}{2}, \]
oder
\[ m = \frac{2n(r-1)}{r} - 3. \]
Erstlich also muss $r$ ein Factor von $2n$ sein, setzt man ferner $m < n$, so findet man $r < 2 + \dfrac{6}{n-3}$. Diesen Bedingungen genügt im Allgemeinen nur $r = 2$; ausserdem aber $r = 4$ für $n = 4$, und $r = 3$ für $n = 6$. Man erhält daher folgende Berührungsaufgaben, welche als vollständig bestimmt von besonderem Interesse sind:

\emph{Die Curve $n^{\text{ter}}$ Ordnung soll von einer Curve der $(n-3)^{\text{ten}}$ Ordnung in $\dfrac{n.n-3}{2}$ Punkten zweipunktig berührt werden;}
wovon die Aufgabe der Doppeltangenten bei den Curven vierter Ordnung ein besonderer Fall ist;

\emph{eine gegebene Curve vierter Ordnung soll von einer Curve dritter Ordnung in drei Punkten vierpunktig berührt werden;}

\origpage{205}
\emph{eine gegebene Curve sechster Ordnung soll von einer Curve fünfter Ordnung in zehn Punkten dreipunktig berührt werden.}

Die beiden letzten Aufgaben sind in dem oben angegebenen Satze bereits enthalten; dennoch mögen die allgemeinen Resultate für diese beiden Fälle hier noch näher specialisirt werden.

Für das zweite Problem erhält man folgende Bestimmung:

\emph{Es giebt 4096 Curven dritter Ordnung, welche eine gegebene Curve vierter Ordnung in drei Punkten vierpunktig berühren.}

Die 4096 Curven entsprechen den 4096 verschiedenen Systemen der $A$, welche man für $p = 3$ bilden kann, wenn die Zahlen $m$, $q$ die Werthe 0, 1, 2, 3 erhalten dürfen. Die Berührungspunkte jeder Curve sind durch eine Gleichung dritten Grades bestimmt; die Coefficienten der Gleichung sind $\theta$-Functionen mit den Argumenten:
\[
\left\{
\begin{aligned}
v_{1} &= u_{1}^{(1)} + u_{1}^{(2)} + u_{1}^{(3)} \equiv \frac{A_{1}}{4}, \\
v_{2} &= u_{2}^{(1)} + u_{2}^{(2)} + u_{2}^{(3)} \equiv \frac{A_{2}}{4}, \\
v_{3} &= u_{3}^{(1)} + u_{3}^{(2)} + u_{3}^{(3)} \equiv \frac{A_{3}}{4}.
\end{aligned}
\right.
\tag{10.}
\]
Bezeichnet man die verschiedenen Berührungscurven entsprechenden Zahlen $m$, $q$ durch obere Indices, \emph{so hat man immer die Berührungspunkte dreier Curven auf einer neuen Curve dritter Ordnung}, sobald
\[
\begin{aligned}
m'_{i} + m''_{i} + m'''_{i} + m''''_{i} &\equiv 0 \quad (\text{mod. } 4), \\
q'_{i} + q''_{i} + q'''_{i} + q''''_{i} &\equiv 0 \quad (\text{mod. } 4).
\end{aligned}
\]
Die so erhaltenen Curven dritter Ordnung theilen sich in drei Classen; die Curven der ersten berühren zweipunktig in den Berührungspunkten zweier der obigen Curven; die der zweiten berühren in den Berührungspunkten einer Curve zweipunktig, und gehen durch die Berührungspunkte zweier anderen hindurch; die der dritten Classe endlich gehen durch die Berührungspunkte von vier verschiedenen Curven. Die Anzahl aller dieser Curven ist sehr gross, und man erhält sie leicht durch eine einfache Abzählung, was ich hier übergehen zu können glaube. Weiteres über die Gruppirung dieser Berührungscurven siehe unten.

Soll eine Curve sechster Ordnung von einer Curve fünfter Ordnung in zehn Punkten dreipunktig berührt werden, so hat man $p = 10$, $m = 5$, $n = 6$, $r = 3$; jedes System von zehn Berührungspunkten wird mittelst einer Gleichung

\origpage{206}
zehnten Grades gefunden, deren Coefficienten $\theta$-Functionen sind mit den Argumenten:
\[
\begin{aligned}
v_{1} &= u_{1}^{(1)} + u_{1}^{(2)} + \cdots + u_{1}^{(10)} \equiv \frac{A_{1}}{3}, \\
v_{2} &= u_{2}^{(1)} + u_{2}^{(2)} + \cdots + u_{2}^{(10)} \equiv \frac{A_{2}}{3}, \\
&\ \cdot\quad\cdot\quad\cdot\quad\cdot\quad\cdot\quad\cdot\quad\cdot \\
v_{10} &= u_{10}^{(1)} + u_{10}^{(2)} + \cdots + u_{10}^{(10)} \equiv \frac{A_{10}}{3}.
\end{aligned}
\]
\emph{Die Anzahl der Berührungscurven ist $3^{20}$.} Indem man die Gleichungen dieser Art, welche zu drei verschiedenen Berührungscurven gehören, addirt, und dann
\[ \frac{A_{k} + A'_{k} + A''_{k}}{3} \equiv 0 \]
setzt, findet man den Satz:

\emph{Jede zwei Systeme von Berührungspunkten liegen auf einer Curve fünfter Ordnung, welche die Curve sechster Ordnung noch in den Berührungspunkten einer dritten Berührungscurve schneidet. Solcher Curven fünfter Ordnung giebt es also $\dfrac{3^{20}.3^{20}-1}{6}$.}

\section*{§.~8.}
\subsection*{Curven $(n-3)^{\text{ter}}$ Ordnung, welche eine gegebene Curve $n^{\text{ter}}$ Ordnung in $\dfrac{n.n-3}{2}$ Punkten zweipunktig berühren.}

Die Aufgabe, \emph{eine Curve $(n-3)^{\text{ter}}$ Ordnung so legen, dass sie eine gegebene Curve $n^{\text{ter}}$ Ordnung in $\dfrac{n.n-3}{2}$ Punkten zweipunktig berührt}, ist bei weitem wichtiger als die vorigen, und von hervorragendem Interesse. Indem wir bei der bisher benutzten Bezeichnung bleiben, stellen sich die Gleichungen zwischen den Integralen, auf welche diese Aufgabe führt, folgendermassen dar:
\[
\left\{
\begin{aligned}
u_{1}^{(1)} + u_{1}^{(2)} + \cdots + u_{1}^{\left(\frac{n.n-3}{2}\right)} &\equiv \frac{A_{1}}{2}, \\
u_{2}^{(1)} + u_{2}^{(2)} + \cdots + u_{2}^{\left(\frac{n.n-3}{2}\right)} &\equiv \frac{A_{2}}{2}, \\
&\ \cdot\quad\cdot\quad\cdot\quad\cdot\quad\cdot\quad\cdot\quad\cdot \\
u_{\left(\frac{n-1.n-2}{2}\right)}^{(1)} + u_{\left(\frac{n-1.n-2}{2}\right)}^{(2)} + \cdots + u_{\left(\frac{n-1.n-2}{2}\right)}^{\left(\frac{n.n-3}{2}\right)} &\equiv \frac{A_{\left(\frac{n-1.n-2}{2}\right)}}{2}.
\end{aligned}
\right.
\tag{11.}
\]

\origpage{207}

Die Anzahl dieser Gleichungen ist um eins grösser als die der Unbekannten. Dennoch sind die Gleichungen unter einander verträglich, wie jetzt gezeigt werden soll; es ist aber keine eine Folge der anderen, vielmehr können sie nur für gewisse Systeme der $A$ neben einander bestehen.

Herr \emph{Riemann} hat a.\ a.\ O. p.\ 147 den Satz bewiesen, dass man ein gegebenes System von $p$ Grössen immer und nur dann, und im Allgemeinen nur auf eine Weise, den Summen von je $p-1$ entsprechenden Integralen gleich setzen könne, wenn $\theta$ für die gegebenen Grössen als Argumente verschwindet. In unserem Falle ist $p = \frac{n-1\,.\,n-2}{2}$, $p-1 = \frac{n\,.\,n-3}{2}$; es giebt also stets eine Berührungscurve der gesuchten Art, sobald $\theta$, mit den Argumenten $\frac{A_{1}}{2}$, $\frac{A_{2}}{2}$, \ldots $\frac{A_{p}}{2}$ gebildet, gleich Null wird. Man muss also haben:
\[
\begin{aligned}
0 &= \mathop{\Sigma}\limits_{n=-\infty}^{n=+\infty} e^{\mathop{\Sigma}\limits_{h,k=1}^{h,k=p} a_{hk}\,n_{h}\,(n_{k}+q_{k}) \,+\, i\pi \mathop{\Sigma}\limits_{k=1}^{k=p} n_{k}m_{k}}\\[2ex]
&= e^{-\frac{i\pi}{2}\Sigma m_{k}q_{k} \,-\, \frac{1}{4}\Sigma\Sigma a_{hk}\,q_{h}q_{k}}.\Sigma e^{\Sigma\Sigma \frac{2n_{h}+q_{h}}{2}.\frac{2n_{k}+q_{k}}{2}\,a_{hk} \,+\, i\pi\Sigma \frac{2n_{k}+q_{k}}{2}\,m_{k}}.
\end{aligned}
\]
Die in der letzten Form auftretende Summe geht, wenn man $-(n_{k}+q_{k})$, $-(n_{h}+q_{h})$ für $n_{k}$, $n_{h}$ setzt, auch über in
\[
\Sigma e^{\Sigma\Sigma \frac{2n_{h}+q_{h}}{2}.\frac{2n_{k}+q_{k}}{2}\,a_{hk} \,-\, i\pi \Sigma \frac{2n_{k}+q_{k}}{2}\,m_{k}},
\]
und wenn man diese Summe zu der vorigen, mit welcher sie identisch ist, addirt, findet sich:
\[
0 = \left\{ \Sigma e^{\Sigma\Sigma \frac{2n_{h}+q_{h}}{2}.\frac{2n_{k}+q_{k}}{2}\,a_{hk}}.(-1)^{\Sigma n_{k}m_{k}} \right\} \cos\left(\frac{\Sigma q_{k}m_{k}}{2}\pi\right).
\]
Die betrachtete $\theta$-Function verschwindet also immer, wenn
\[
q_{1}m_{1}+q_{2}m_{2}+\cdots+q_{p}m_{p} \;=\; 2h+1, \tag{12.}
\]
und dieses ist die Bedingung für das Zusammenbestehen der Gleichungen (11.).

Untersuchen wir nun die Anzahl der Combinationen, durch welche die Gleichungen (12.) erfüllt werden kann. Seien $k$ der Grössen $q$ gleich 0, die übrigen $(p-k)$ gleich 1; die den ersten Grössen $q$ entsprechenden $m$ können dann 0 oder 1 sein, was $2^{k}$ Combinationen giebt; die Summe der anderen $m$ aber muss ungerade sein, oder es ist immer das letzte durch die übrigen $p-k-1$

\origpage{208}

bestimmt. Die letzteren $m$ können daher noch auf $2^{p-k-1}$ Weisen gewählt werden, oder die $m$ im Ganzen auf $2^{p-1}$ Arten. Da die $q$ auf $\frac{p.p-1\ldots p-k+1}{1\,.\,2\ldots k}$ Arten den angegebenen Bedingungen genügen können, so giebt es
\[
\frac{p.p-1\ldots p-k+1}{1.2\ldots k}.2^{p-1}
\]
Systeme der $A$, bei denen $k$ der Grössen $q$ gleich 0 sind. Setzen wir nun $k = 0,\ 1,\ 2,\ \ldots p-1$, so finden sich im Ganzen
\[
2^{p-1}\left(1+\frac{p}{1}+\frac{p.p-1}{1\,.\,2}+\cdots+p\right) \;=\; 2^{p-1}(2^{p}-1)
\]
Systeme der $A$, oder der Satz:

\emph{Es giebt $2^{\frac{n\,.\,n-3}{2}}\left(2^{\frac{n-1\,.\,n-2}{2}}-1\right)$ Curven $(n-3)^{\text{ter}}$ Ordnung, welche eine gegebene Curve $n^{\text{ter}}$ Ordnung in $\frac{n\,.\,n-3}{2}$ Punkten zweipunktig berühren.}

Für $n=4$ giebt dies die 28 Doppeltangenten.

Die Berührungspunkte jeder Curve können mittelst einer Gleichung $(p-1)^{\text{ten}}$ Grades gefunden werden. Man braucht nämlich nur auf beiden Seiten der Gleichungen (11.) die einem beliebig gegebenen Punkte entsprechenden $u$ zu addiren, und man hat dann die Argumente von $\theta$-Functionen vor sich. Aus diesen setzen sich die Coefficienten einer Gleichung $p^{\text{ten}}$ Grades zusammen, von der eine Wurzel (dem beliebig gewählten Punkte entsprechend) bekannt ist.

Die Berührungspunkte von $\mu$ Curven dieser Art liegen jedesmal auf einer Curve der Ordnung $\frac{\mu(n-3)}{2}$, sobald das System besteht:
\[
\frac{A_{k}^{(1)}+A_{k}^{(2)}+\cdots+A_{k}^{(\mu)}}{2} \;\equiv\; 0,
\]
bei welchem die oberen Indices sich auf die verschiedenen betrachteten Curven beziehen. Man überblickt das Feld von Sätzen und Abzählungen, welche sich hieraus ergeben, und welche die natürliche Verallgemeinerung derjenigen Resultate bilden, auf welche man in der Theorie der Doppeltangenten der Curven vierter Ordnung geführt ist. Ich bemerke nur noch einen Satz, welcher ebenfalls in speciellster Fassung bei den Curven vierter Ordnung bereits bekannt ist. Betrachtet man nämlich die Curven $2h-(n-3)^{\text{ter}}$ Ordnung, welche eine Curve $n^{\text{ter}}$ Ordnung in $hn-\frac{n\,.\,n-3}{2}$ Punkten zweipunktig berühren, und addirt zu den entsprechenden Gleichungen
\[
u_{k}^{(1)}+u_{k}^{(2)}+\cdots+u_{k}^{\left(hn-\frac{n\,.\,n-3}{2}\right)} \;\equiv\; \frac{A_{k}}{2}
\]

\origpage{209}

die Gleichungen (11.) unter Voraussetzung desselben Systems der $A$, so findet sich, dass die ersten $hn-\frac{n\,.\,n-3}{2}$ Berührungspunkte mit den $\frac{n\,.\,n-3}{2}$ Berührungspunkten der Curve $(n-3)^{\text{ter}}$ Ordnung in einer Curve $h^{\text{ter}}$ Ordnung liegen. Nun ist die Anzahl der Systeme von Berührungscurven $2h-(n-3)^{\text{ter}}$ Ordnung gleich $2^{2p}$, und nur, wenn $n$ ungerade ist, muss man das System ausschliessen, für welches das System der $A$ verschwindet. Man hat also folgenden Satz:

\emph{Unter den $2^{2p}\left(p = \frac{n-1\,.\,n-2}{2}\right)$ Systemen von Berührungscurven, welche bei geradem $n$ von ungerader Ordnung $m$, oder bei ungeradem $n$ von gerader Ordnung $m$ sind, und welche die gegebene Curve $n^{\text{ter}}$ Ordnung überall, wo sie derselben begegnen, zweipunktig berühren, giebt es immer $2^{p-1}(2^{p}-1)$ solche, dass die Berührungspunkte jeder Curve des Systems mit denen einer bestimmten der obigen Berührungscurven $(n-3)^{\text{ter}}$ Ordnung in einer Curve $\frac{m+n-3}{2}{}^{\text{ter}}$ Ordnung liegen; und es giebt bei geradem $n$ $2^{p-1}(2^{p}+1)$, bei ungeradem $2^{p-1}(2^{p}+1)-1$ solcher Systeme, bei denen die Berührungspunkte niemals auf einer Curve $\frac{m+n-3}{2}{}^{\text{ter}}$ Ordnung liegen.}

\section*{§.~9.}

\subsection*{Die Doppeltangenten der Curven vierter Ordnung.}

Die entwickelten allgemeinen Resultate sollen jetzt für die Curven vierter Ordnung genauer ausgeführt werden.

Die Anzahl der Doppeltangenten ist nach §.~8 gleich 28, wie bekannt. Man kann aber nach dem Vorigen sie schematisch darstellen, wie folgt. Jede Doppeltangente mag bezeichnet werden durch die in Klammern geschlossene Reihe der ihr entsprechenden Zahlen $m$, $q$, also durch:
\[
(m_{1},\, m_{2},\, m_{3};\, q_{1},\, q_{2},\, q_{3}).
\]
Die Zahlen können 0 oder 1 sein, doch so, dass immer
\[
m_{1}q_{1}+m_{2}q_{2}+m_{3}q_{3} \;\equiv\; 1 \quad \text{(mod.\ 2)}.
\]
Nehmen wir also für die $m$ alle möglichen aus den Zahlen 0 und 1 bestehenden Systeme, $0,0,0$ ausgeschlossen, und bestimmen obiger Gleichung gemäss die zugehörigen Systeme der $q$, so erhalten wir die 28 Doppeltangenten in folgendem Schema:
\[
\begin{array}{ccccccc}
(100,\,100) & (010,\,010) & (001,\,001) & (011,\,010) & (101,\,100) & (110,\,100) & (111,\,100)\\
(100,\,110) & (010,\,110) & (001,\,011) & (011,\,001) & (101,\,110) & (110,\,101) & (111,\,010)\\
(100,\,101) & (010,\,011) & (001,\,101) & (011,\,110) & (101,\,001) & (110,\,010) & (111,\,001)\\
(100,\,111) & (010,\,111) & (001,\,111) & (011,\,101) & (101,\,011) & (110,\,011) & (111,\,111).
\end{array}
\]

\origpage{210}

Es giebt nach §.~5 63 Systeme von Kegelschnitten (\emph{Hesse}, dieses Journal Band 49.), welche die Curve in 4 Punkten berühren, entsprechend den 63 möglichen Combinationen
\[
(m_{1},\, m_{2},\, m_{3};\, q_{1},\, q_{2},\, q_{3}),
\]
die Combination $(000,\, 000)$ allein ausgeschlossen.

Ein solcher Berührungskegelschnitt zerfällt immer in zwei Doppeltangenten
\[
(m'_{1},\, m'_{2},\, m'_{3};\; q'_{1},\, q'_{2},\, q'_{3}), \qquad (m''_{1},\, m''_{2},\, m''_{3};\; q''_{1},\, q''_{2},\, q''_{3}),
\]
sobald
\[
m'_{i}+m''_{i} \;\equiv\; m_{i}, \qquad q'_{i}+q''_{i} \;\equiv\; q_{i} \quad \text{(mod.\ 2)}, \tag{13.}
\]
wobei ausserdem sein muss:
\[
\left\{
\begin{aligned}
m'_{1}q'_{1}+m'_{2}q'_{2}+m'_{3}q'_{3} &\;\equiv\; 1,\\
m''_{1}q''_{1}+m''_{2}q''_{2}+m''_{3}q''_{3} &\;\equiv\; 1.
\end{aligned}
\right. \tag{14.}
\]
Es sind dabei entweder die $q$ oder die $m$ nicht sämmtlich Null; nehmen wir an, dies trete bei den $q$ ein; sonst wären in der folgenden Bestimmung nur die $q$ mit den $m$ zu vertauschen. Dann zeigen erstlich die Gleichungen (13.), dass man für gegebene $q_{i}$ die $q'_{i}$, $q''_{i}$ auf sechs verschiedene Arten wählen kann. Ueberhaupt nämlich existiren acht Combinationen der $q'$, $q''$, welche jene Gleichungen erfüllen; von diesen aber sind die beiden auszuschliessen, für welche sämmtliche $q'$ oder sämmtliche $q''$ Null sind (was nicht zusammen eintreten kann). Diese sechs Arten geben drei Paare von Systemen $q'$, $q''$, indem immer zwei Arten sich nur durch Vertauschung der $q'$ mit den $q''$ von einander unterscheiden.

In jedem Paar von Systemen ist wenigstens ein Paar entsprechender $q$ ungleich, also 0 und 1, demnach irgend ein anderes 1 und 0 oder 1 und 1. Es existiren daher ein $q'$ und ein $q''$ mit verschiedenem unteren Index, welche gleich 1 sind. Bestimmt man nun die dem dritten Index entsprechenden $m$ so, dass sie der betreffenden Gleichung (13.) genügen (was auf zwei Arten geschehen kann), so bestimmen die übrigen $m$ sich aus (14.) und (13.) successive. Jedem der drei Paare $q$ entsprechen also zwei Paare $m$; oder man hat den Satz:

\emph{In jedem System von Berührungskegelschnitten kommen sechs Paare von Doppeltangenten vor.} (\emph{Hesse}, l.\ c., \emph{Salmon}, higher plane curves.)

Da für irgend zwei solcher demselben System angehöriger Paare die Summe aller $m_{i}$ so wie die Summe aller $q_{i}$ stets gerade ist, so hat man ferner den Satz:

\origpage{211}

\emph{Die Berührungspunkte je zweier Paare, welche demselben System angehören, liegen auf einem Kegelschnitt.} (ib., ib.)

Die Zahl der so erhaltenen Kegelschnitte ist bekanntlich gleich der der 63 Systeme, multiplicirt mit der Anzahl 15 der Combinationen von 6 Paaren zu zweien, und dividirt durch 3, da jede vier Tangenten, deren Berührungspunkte auf einem Kegelschnitt liegen, auf drei Arten in zwei Paare zerlegt werden können, also jeder Kegelschnitt dreimal vorkommt. So ist die Zahl dieser Kegelschnitte 315. Das obige Schema erlaubt sofort, die solchen Kegelschnitten entsprechenden Doppeltangenten zusammenzustellen.

Der allgemeinere Satz, \emph{dass die Berührungspunkte je zweier Kegelschnitte desselben Systems auf einem Kegelschnitt liegen} (\emph{Hesse}, l.\ c.), folgt aus §.~6.

Es giebt ferner nach §.~6 \emph{64 Systeme von Curven dritter Ordnung, welche die Curve vierter Ordnung in 6 Punkten berühren.} Die den Berührungspunkten entsprechenden Integrale genügen dann immer den Gleichungen:
\[
u_k^{(1)}+u_k^{(2)}+\cdots+u_k^{(6)} \equiv \frac{A_k}{2}. \tag{15.}
\]
Man übersieht nun, dass diese 64 Systeme sich sofort in zwei Classen sondern, je nachdem das System der $A$ zugleich bei den Doppeltangenten vorkommt oder nicht; je nachdem also $\Sigma mq \equiv 1$ (was 28 mal geschieht) oder $\Sigma mq \equiv 0$ (was auf 36 Arten geschehen kann). Hat man nämlich für irgend eine Doppeltangente
\[
u_k^{(7)}+u_k^{(8)} \equiv \frac{A_k}{2}, \tag{16.}
\]
so folgt
\[
u_k^{(1)}+u_k^{(2)}+\cdots+u_k^{(8)} \equiv 0,
\]
d.\ h.\ die Berührungspunkte der Curve dritter Ordnung liegen mit denen der Doppeltangente auf einem Kegelschnitt; während, wenn es für ein System der $A$ keine Gleichung wie (16.) giebt, auch die Berührungspunkte der Curve auf keinem Kegelschnitt liegen können, da die Gleichung (15.) dann eine passende Ergänzung nicht zulässt. Man hat also den Satz (\emph{Hesse}, l.\ c.):

\emph{Von den 64 Systemen der Curven dritter Ordnung haben 28 die Eigenschaft, dass die 6 Berührungspunkte einer Curve auf einem Kegelschnitt liegen, und jeder Kegelschnitt geht dann noch immer durch die Berührungspunkte einer Doppeltangente, welche für dasselbe System constant bleibt. Die Berührungspunkte bei Curven der übrigen 36 Systeme liegen nicht in einem Kegelschnitt.}

\origpage{212}

Dagegen haben nach §.~6 beide Classen die gemeinsame Eigenschaft, \emph{dass die Berührungspunkte je zweier Curven desselben Systems auf einer Curve dritter Ordnung liegen.} (\emph{Hesse}, l.\ c.)

Zu diesen Systemen gehören auch die 4096 Curven dritter Ordnung, welche die Curve vierter Ordnung in drei Punkten vierpunktig berühren (§.~7). In der That folgt aus (10.)
\[
2\,(u_k^{(1)}+u_k^{(2)}+u_k^{(3)}) \equiv \frac{A_k}{2};
\]
jede Curve gehört also zu demjenigen System von Berührungscurven, dem das System der $A_k$ zukommt. Umgekehrt aber folgt aus dieser Gleichung
\[
u_k^{(1)}+u_k^{(2)}+u_k^{(3)} \equiv \frac{A_k}{4}+\frac{A'_k}{2},
\]
wo in $A_k$, $A'_k$ die Zahlen $m$, $q$ jetzt überall nur die Werthe 0 und 1 erhalten dürfen. Ist über die $A_k$ verfügt, so können die $A'_k$ noch auf 64 Arten gewählt werden; \emph{daher gehören in jedes der 64 Systeme von Berührungscurven dritter Ordnung 64 solche Curven, welche die Curve vierter Ordnung in drei Punkten vierpunktig berühren.} Dies stimmt damit überein, dass nach §.~4 durch jede zwei Punkte einer Curve vierter Ordnung 64 Kegelschnitte gelegt werden können, welche dieselbe in drei Punkten berühren. Legt man diese Kegelschnitte durch die Berührungspunkte einer Doppeltangente, so sind die Berührungspunkte des Kegelschnitts zugleich diejenigen, in welchen eine Curve dritter Ordnung vierpunktig berühren kann. \emph{So erhält man die $28.64$ Berührungscurven, welche den erwähnten 28 Systemen angehören; aber man sieht, dass es noch $36.64$ solcher Berührungscurven giebt, in deren Berührungspunkten nicht zugleich ein Kegelschnitt die Curve vierter Ordnung berühren kann.}

\emph{Ueberhaupt giebt es nach §.~5 je 63 Systeme von Curven $(2m)^{\text{ter}}$ Ordnung und je 64 Systeme von Curven $(2m+1)^{\text{ter}}$ Ordnung, welche die Curve vierter Ordnung zweipunktig berühren, wo sie derselben begegnen. Und diese letzteren 64 Systeme theilen sich immer in 28 und 36 so, dass die Berührungspunkte einer Curve aus einem der 28 Systeme immer mit den Berührungspunkten einer der 28 Doppeltangenten auf einer Curve $(m+1)^{\text{ter}}$ Ordnung liegen, während etwas Aehnliches bei den anderen 36 Systemen niemals stattfindet. Aber alle Systeme haben die Eigenschaft, dass die Berührungspunkte je zweier Curven desselben Systems auf einer Curve desselben Grades liegen.}

Man kann von diesen Sätzen leicht zu einer genaueren Discussion derjenigen Fälle übergehen, in denen die Berührungspunkte von 6, 8 etc. Doppeltangenten

\origpage{213}

auf einer Curve dritter, vierter etc. Ordnung liegen, worüber von \emph{Hesse} und \emph{Steiner} (Bd.\ 49 dieses Journals) Sätze gefunden sind. Ich übergehe diese Ausführungen, welche nach den entwickelten Principien durch blosse Abzählungen erhalten werden können.

Von Curven, welche die Curve vierter Ordnung mehrpunktig berühren, will ich wenigstens noch die der dritten Ordnung anführen. Aus §.~6 folgt der Satz:

\emph{Es giebt 728 Systeme von Curven dritter Ordnung, welche eine Curve vierter Ordnung in vier Punkten dreipunktig berühren.}

Zwischen den Integralen bestehen hier die Gleichungen:
\[
u_k^{(1)}+u_k^{(2)}+u_k^{(3)}+u_k^{(4)} \equiv \frac{A_k}{3};
\]
das System der $A$ darf niemals verschwinden. Betrachtet man eine Curve, einem anderen System angehörig, für welches immer $2A_k$ an die Stelle von $A_k$ tritt, und addirt die einer solchen Curve entsprechenden Gleichungen zu den vorigen, so findet sich die Summe analoger Integrale immer der Null congruent. Man hat daher folgenden Satz, welcher (unter anderen) das Verhalten dieser Systeme charakterisirt:

\emph{Diese 728 Systeme theilen sich in 364 Paare zu zweien so, dass die Berührungspunkte jeder Curve des einen Systems mit denen jeder Curve des anderen Systems auf einem Kegelschnitt liegen.}

Aehnliche Sätze kann man für die Berührungscurven dieser und anderer Ordnungen in Menge aufstellen.

\section*{§.~10.}

\subsection*{Ueber einige Sätze von Steiner, die Curven vierter Ordnung betreffend.}

Ich benutze diese Gelegenheit, um einige Sätze über Curven vierter Ordnung zu beweisen, welche \emph{Steiner} (dieses Journal Bd.\ 49, p.\ 266 folgg.) angegeben hat, obwohl dieselben nicht in die Reihe der hier angestellten Betrachtungen gehören.

Nach \emph{Hesse} (dieses Journal Band 49, p.\ 261) kommt die Aufstellung eines Berührungskegelschnitts darauf hinaus, der Gleichung der Curve die Form
\[
uw - v^{2} = 0
\]
zu geben, wo $u$, $v$, $w$ Ausdrücke der zweiten Ordnung sind. Alle Kegelschnitte, welche demselben der 63 Systeme angehören, sind dann in der Form enthalten
\[
u + 2\lambda v + \lambda^{2} w = 0,
\]

\origpage{214}

und man findet die 6 in dem System vorkommenden Tangentenpaare durch eine Gleichung sechsten Grades in $\lambda$, indem man die Determinante dieses Kegelschnitts verschwinden lässt. Für den Schnittpunkt $x$ der Tangenten eines Paars müssen dann die Differentialquotienten
\[
u_i + 2\lambda v_i + \lambda^{2} w_i
\]
verschwinden, und man findet nach einem bekannten Verfahren die Producte und Quadrate der Coordinaten $x$ den Unterdeterminanten des Kegelschnittes $u+2\lambda v+\lambda^{2}w = 0$ proportional. Diese Ausdrücke haben, wie man sofort sieht, die Form:
\[
\varrho\, x_i x_k = p_{ik}+\lambda q_{ik}+\lambda^{2} r_{ik}+\lambda^{3} s_{ik}+\lambda^{4} t_{ik},
\]
wo $\varrho$ einen willkürlichen Factor bedeutet, und wo die $p$, $q$, $\ldots$ nur noch von den Coefficienten der Curve abhängen. Bildet man daher die Determinante dieser sechs Gleichungen, indem man $\varrho$ und $\lambda$ eliminirt, so hat man eine Gleichung zweiten Grades für die $x$, welche von der gewählten Wurzel $\lambda$ ganz unabhängig ist. Und so hat man den von \emph{Steiner} a.\ a.\ O.\ gegebenen Satz:

\emph{Die Doppelpunkte der 6 demselben System angehörigen Paare von Doppeltangenten liegen auf einem Kegelschnitt.}

Da identisch
\[
(uw-v^{2})(k-\lambda)^{2} = (u+2kv+k^{2}w)(u+2\lambda v+\lambda^{2}w)-(u+(k+\lambda)v+k\lambda w)^{2},
\]
so ist
\[
u+(k+\lambda)v+k\lambda w = 0,
\]
oder was dasselbe ist, jeder Kegelschnitt von der Form
\[
\alpha u+\beta v+\gamma w = 0
\]
ist ein solcher, welcher durch die Berührungspunkte zweier Kegelschnitte des Systems hindurchgeht. Lässt man denselben in ein Linienpaar $p$, $q$ zerfallen, d.\ h.\ setzt man:
\[
\alpha u_{ik}+\beta v_{ik}+\gamma w_{ik} = p_i q_k+q_i p_k,
\]
so genügen die Coordinaten einer Geraden des Paars der Gleichung, welche durch Elimination von $\alpha$, $\beta$, $\gamma$, $q_1$, $q_2$, $q_3$ entsteht:
\[
K = \begin{vmatrix}
u_{11} & v_{11} & w_{11} & 2p_1 & 0 & 0\\
u_{22} & v_{22} & w_{22} & 0 & 2p_2 & 0\\
u_{33} & v_{33} & w_{33} & 0 & 0 & 2p_3\\
u_{23} & v_{23} & w_{23} & 0 & p_3 & p_2\\
u_{31} & v_{31} & w_{31} & p_3 & 0 & p_1\\
u_{12} & v_{12} & w_{12} & p_2 & p_1 & 0
\end{vmatrix} = 0.
\]

\origpage{215}

Die Doppelpunkte der Paare müssen zugleich den Gleichungen genügen:
\[
\alpha u_i+\beta v_i+\gamma w_i = 0,
\]
oder, indem man $\alpha$, $\beta$, $\gamma$ eliminirt, der Gleichung
\[
G = \begin{vmatrix}
u_1 & v_1 & w_1\\
u_2 & v_2 & w_2\\
u_3 & v_3 & w_3
\end{vmatrix} = 0.
\]
Man hat also den Satz:

\emph{Wenn man die vollständigen Vierecke construirt, deren Ecken die Berührungspunkte von Kegelschnitten desselben Systems sind, so bilden die Durchschnitte gegenüberliegender Seiten die Curve dritter Ordnung $G=0$, und die Seiten selbst umhüllen die Curve dritter Classe $K=0$.}

Insbesondere folgen daraus die von \emph{Steiner} a.\ a.\ O.\ gegebenen Sätze als specielle Fälle:

\emph{Auf der Curve $G=0$ liegen die Doppelpunkte der 6 dem System angehörigen Paare von Doppeltangenten, so wie die Schnittpunkte der Geraden, welche die Berührungspunkte jedes Paars kreuzweise verbinden.}

Und:

\emph{Die Curve $K=0$ wird berührt von den sechs Paaren dem System angehöriger Doppeltangenten, so wie von den Geraden, welche die Berührungspunkte der Paare kreuzweise verbinden.}

Lässt man in einem vollständigen Viereck 1, 2, 3, 4, dessen Ecken Berührungspunkte eines Kegelschnitts sein sollen, zwei, etwa 2 und 3 zusammenfallen, so fallen auch 5, 6 (siehe die Fig.) in denselben Punkt. Es gehören aber die Punkte 5, 6 der Curve $G=0$ an; diese schneidet also in dem betreffenden Punkte die gegebene Curve vierter Ordnung, indem erstere die Gerade 5, 6, letztere die Gerade 2, 3 zur Tangente hat. Die Schnittpunkte von $G=0$ mit der gegebenen Curve haben also die Eigenschaft, dass in ihnen ein Kegelschnitt die Curve vierpunktig berühren kann, der sie ausserdem noch (in 1 und 4) zweimal zweipunktig berührt. In jedem System giebt es demnach 12 Kegelschnitte dieser Art.

\rmfigure{figures/fig-215.png}{Fig.}{Line diagram of a degenerate complete quadrilateral: points 5 and 2, 3 near the top, with straight lines through them meeting at points labelled 6, 8, 4 in the interior and at 1 (lower left) and 7 (right), illustrating the harmonic system 1, 8, 4, 7.}

Da ferner die Punkte 1, 8, 4, 7 ein harmonisches System bilden, so bilden auch in den betrachteten Punkten die Tangente an die gegebene Curve, die Tangente an $G=0$, und die von der vierpunktigen Berührungsstelle nach

\origpage{216}

den zweipunktigen gezogenen Geraden ein harmonisches System. --- Der Punkt 7 gehört der Curve $G=0$ an; er wird hier der Durchschnitt der Tangente an der vierpunktigen Berührungsstelle mit der Sehne, welche die zweipunktigen Berührungsstellen verbindet.

Die Tangente 2, 3 an einer der vierpunktigen Berührungsstellen ist Tangente an $K=0$; ebenso die Geraden 2, 1; 3, 4, welche jene Stelle mit denen der zweipunktigen Berührung verbinden; da nun auch 3, 1; 2, 4 Tangenten dieser Curve sind, und diese mit den ebengenannten Geraden unendlich nahe zusammenfallen, so müssen 1, 4, die Stellen der zweipunktigen Berührung, als Schnittpunkte nächster Tangenten auf $K=0$ liegen. Diese Curve ist von der sechsten Ordnung; die 24 in einem System vorkommenden Punkte 1, 4 stellen also das vollständige Schnittsystem der Curve $K=0$ mit der gegebenen Curve vor. So sind die folgenden von \emph{Steiner} gegebenen Sätze bewiesen:

\emph{Es giebt $63.12 = 756$ Kegelschnitte, welche eine gegebene Curve vierter Ordnung in einem Punkte $a$ vierpunktig, und ausserdem noch in zwei anderen, $b$, $c$, zweipunktig berühren.}

\emph{Die Punkte $a$ sind in jedem System die 12 Durchschnitte von $G=0$ mit der gegebenen Curve; auf $G=0$ liegen auch $b$, $c$ und der Schnittpunkt der Tangente der gegebenen Curve im Punkte $a$ mit der Sehne $b$, $c$.}

\emph{In jedem Punkte $a$ bilden die nach $b$, $c$} \uncertain{gezogenen} \emph{Geraden mit den Tangenten der gegebenen Curve und der Curve $G=0$ ein harmonisches System.}

\emph{Die Tangente in $a$, sowie die drei} \uncertain{Geraden} \emph{$ab$, $ac$, $bc$ berühren die Curve $K=0$, und zwar $ab$, $ac$ in $b$ und} \uncertain{$c$} \emph{selbst, so dass die Punkte $b$, $c$ in jedem System den vollständigen Durchschnitt von $K=0$ mit der gegebenen Curve darstellen.}

\emph{Steiner} entwickelt endlich Zusammenhänge zwischen den Curven $G=0$, $K=0$. Die Beziehung, in welcher diese Curven zu einander stehen, kann man folgendermassen erkennen.

Betrachten wir statt der Functionen $u$, $v$, $w$ irgend welche lineare Verbindungen derselben
\[
\begin{aligned}
U_1 &= a_1 u+b_1 v+c_1 w,\\
U_2 &= a_2 u+b_2 v+c_2 w,\\
U_3 &= a_3 u+b_3 v+c_3 w.
\end{aligned}
\]

\origpage{217}

Diese drei neuen Functionen können immer so bestimmt werden, dass sie die partiellen Differentialquotienten einer Function dritter Ordnung nach $x_1$, $x_2$, $x_3$ sind. Hierzu ist nöthig, dass folgende identische Gleichungen bestehen:
\[
\begin{aligned}
a_2 u_3+b_2 v_3+c_2 w_3 &= a_3 u_2+b_3 v_2+c_3 w_2,\\
a_3 u_1+b_3 v_1+c_3 w_1 &= a_1 u_3+b_1 v_3+c_1 w_3,\\
a_1 u_2+b_1 v_2+c_1 w_2 &= a_2 u_1+b_2 v_1+c_2 w_1.
\end{aligned}
\]
Diese Gleichungen sind für die $x$ linear; setzt man also die Coefficienten auf beiden Seiten einander gleich, so findet man 9 lineare Gleichungen, welche in Bezug auf die gesuchten Grössen $a$, $b$, $c$ homogen sind. Dennoch sind sie erfüllbar; denn man bemerkt sofort, dass ihre Determinante eine überschlagene ist, und demnach, weil von ungeradem Grade, identisch verschwindet.

Nun ändern sich offenbar die Ausdrücke $G$ und $K$ keineswegs, wenn man bei der Bildung derselben statt $u$, $v$, $w$ irgend welche lineare Verbindungen derselben zu Grunde legt, z.\ B.\ die eben gefundenen partiellen Differentialquotienten einer homogenen Function $f$ dritter Ordnung. Dann aber hat man, bis auf einen Zahlenfactor:
\[
G = \Delta_f, \qquad K = S_f,
\]
wo $\Delta_f$ die Covariante dritter Ordnung, $S_f$ die erste zugehörige Form dritter Classe von $f$ bedeutet. Aus der bekannten Rolle, welche die Curven $\Delta_f=0$, $S_f=0$ in der Theorie der Curven dritter Ordnung spielen, folgen daher unmittelbar die von \emph{Steiner} angegebenen Beziehungen. Die Curven $\Delta_f=0$, $S_f=0$ sind dieselben, welche Herr \emph{Cremona} in seiner „Introduzione ad una teoria geometrica della curve piane" als \emph{Hesse}sche und \emph{Cayley}sche Curve bezeichnet hat. Von diesen gelten folgende Sätze:

\emph{Die zu den Wendepunkten der Hesseschen Curve gehörigen harmonischen Polaren sind die Rückkehrtangenten der Cayleyschen Curve.} (\emph{Hesse}, dieses Journal Bd.\ 38, p.\ 252.)

\emph{Die zu den Rückkehrtangenten der Cayleyschen Curve gehörigen harmonischen Punkte sind die Wendepunkte der Hesseschen Curve.} (ib.)

\emph{Die Wendetangenten von $u=0$ berühren die Hessesche und die Cayleysche Curve in denselben 9 Punkten.} (\emph{Cremona}, l.\ c.\ p.\ 116.)

\emph{Jede Tangente der Cayleyschen Curve schneidet die Hessesche Curve in zwei Punkten, deren Tangenten sich auf der Hesseschen Curve treffen, und in einem dritten Punkt, der mit dem Berührungspunkt auf der Cayleyschen Curve und mit den anderen beiden Schnittpunkten ein harmonisches System bildet.} (ib.)

\origpage{218}

\emph{Von jedem Punkte der Hesseschen Curve gehen an die Cayleysche drei Tangenten; die Verbindungslinie der Berührungspunkte zweier ist Tangente der Cayleyschen Curve; die dritte bildet mit der im Ausgangspunkt an die Hessesche Curve gezogenen Tangente und mit den ersten beiden Tangenten ein harmonisches System.} (ib.)

Die blosse Anwendung dieser Sätze giebt die folgenden, von \emph{Steiner} gegebenen Resultate:

\emph{Aus jedem Wendepunkte $w$ der Curve $G=0$ gehen drei Tangenten $Q$, $Q_1$, $Q_2$ an die Curve, welche in $q$, $q_1$, $q_2$ berühren. Die letzteren drei Punkte liegen auf einer Rückkehrtangente $R$ der Curve $K=0$. Diese schneidet $K=0$ ausser $q$ und ausser dem Rückkehrpunkt $r$ in zwei Punkten $q'$, $q''$, deren Tangenten $Q'$, $Q''$ sich ebenfalls in $w$ schneiden. Die Punkte $r$, $q$, $q_1$, $q_2$ bilden ein harmonisches System, die Tangenten $Q$, $Q'$, $Q''$ mit der in $w$ gezogenen Wendetangente ein zweites.}

\section*{§.~11.}

\subsection*{Ueber den Zusammenhang der Curven doppelter Krümmung mit Abelschen Functionen.}

Im Vorigen sind nur ebene Curven behandelt worden. Ich werde jetzt zeigen, wie ähnliche Principien sich für die Behandlung der Curven doppelter Krümmung aufstellen lassen.

Betrachten wir den Durchschnitt zweier algebraischen Flächen. Die eine derselben, $u=0$, sei $m^{\text{ter}}$ Ordnung, die andere, $v=0$, $n^{\text{ter}}$ Ordnung. Die Durchschnittscurve können wir mit zwei Ebenenbüscheln, deren Axen beliebig im Raume liegen, so in Zusammenhang bringen, dass zwei entsprechende Ebenen beider Büschel sich auf der Curve durchschneiden. Sind nämlich die Gleichungen der Büschel
\[
\left\{
\begin{aligned}
(a_1 x_1+a_2 x_2+a_3 x_3+a_4 x_4)+z\,(b_1 x_1+b_2 x_2+b_3 x_3+b_4 x_4) &= 0,\\
(\alpha_1 x_1+\alpha_2 x_2+\alpha_3 x_3+\alpha_4 x_4)+s\,(\beta_1 x_1+\beta_2 x_2+\beta_3 x_3+\beta_4 x_4) &= 0,
\end{aligned}
\right. \tag{17.}
\]
so führt die Elimination der $x$ aus diesen Gleichungen und aus $u=0$, $v=0$ auf eine Gleichung
\[
F(s,z) = 0, \tag{18.}
\]
welche das gegenseitige Abhängigkeitsverhältniss der Parameter der Büschel ausdrückt.

\origpage{219}

Diese Gleichung ist sowohl für $s$ als für $z$ vom Grade $m+n$. Aber sie ist nicht unter allen Umständen irreductibel, sondern zerfällt, wenn die Oberflächen $u=0$, $v=0$ besondere Lagen gegen einander haben, von selbst in Factoren. Sei nämlich ein Theil der Schnittcurve zugleich als Theil des Durchschnitts zweier Flächen $u'=0$, $v'=0$ darstellbar, welche von niedrigerem Grade als $u=0$, $v=0$ sind; sei der zugefügte Theil $a'$ mit einem andern $a''$ zusammen die Schnittcurve zweier Flächen \uncertain{$u''=0$}, $v''=0$, welche wieder von niedrigerem Grade sind, u.\ s.\ w., bis endlich ein Curventheil $a^{(\mu-1)}$ mit einer Curve $a^{(\mu)}$ zusammen der Durchschnitt zweier Flächen $u^{(\mu)}=0$, $v^{(\mu)}=0$ ist, und $a^{(\mu)}$ sich als der vollständige Schnitt zweier Flächen von niedrigerem Grade darstellen lässt. So giebt es unzählige Oberflächenpaare deren Durchschnitt in eine andere Curve und in eine Raumcurve dritter Ordnung zerfällt; diese ist zusammen mit einer geraden Linie als Durchschnitt zweier Flächen zweiter Ordnung darstellbar, und die hinzugefügte Gerade endlich bildet den vollständigen Durchschnitt zweier Ebenen. Denken wir uns nun für alle diese Schnittcurven der Reihe nach die der Gleichung (18.) entsprechenden Gleichungen gebildet;
\[
F(s,z)=0, \qquad F^{(1)}(s,z)=0, \qquad \ldots \qquad F^{(\mu+1)}(s,z)=0,
\]
so müssen folgende Gleichungen bestehen, in welchen die $\varphi$ ganze Functionen bedeuten:
\[
\begin{aligned}
F^{(\mu)}(s,z) &= \varphi^{(\mu)}(s,z).F^{(\mu+1)}(s,z),\\
F^{(\mu-1)}(s,z) &= \varphi^{(\mu-1)}(s,z).\varphi^{(\mu)}(s,z),\\
F^{(\mu-2)}(s,z) &= \varphi^{(\mu-2)}(s,z).\varphi^{(\mu-1)}(s,z),
\end{aligned}
\]
\[
\cdot\quad\cdot\quad\cdot\quad\cdot\quad\cdot\quad\cdot\quad\cdot\quad\cdot\quad\cdot\quad\cdot\quad\cdot\quad\cdot
\]
\[
F(s,z) = \varphi(s,z).\varphi^{(1)}(s,z).
\]
In diesem Falle ist also die Gleichung $F(s,z)=0$ reductibel, und man kann sie als dass Product zweier, oder allgemein zu reden, mehrerer irreductibler Factoren auffassen. Jeder Zweig einer Raumcurve, welcher einem solchen Factor entspricht, kann als \emph{algebraisch definirte einfache Raumcurve} gelten. Ihre Ordnung ist die Ordnung von $s$ und $z$ in dem entsprechenden Factor; Ordnungen, welche immer gleich sein werden, wie aus dem eben angeführten Schema folgt, da nach dem früheren $s$ und $z$ in allen Functionen $F$ zu gleichen Ordnungen vorkommen.

Jede solche einfache Raumcurve begründet nun eine Classe von \emph{Abel}schen Functionen. Sei $k$ die Ordnung der Curve, $\varphi(s,z)=0$ die ihr entsprechende Gleichung.

\origpage{220}

Die Zahl $p$, welche die zugehörigen \emph{Abel}schen Functionen charakterisirt, ist dann nach §.~1.\ gleich $\tfrac{1}{2}\mathrm{w}-(k-1)$, wo $\mathrm{w}$ die Anzahl derjenigen Werthsysteme $s$, $z$ bedeutet, für welche $\varphi'(s)=0$, ohne dass $\varphi'(z)=0$.

Ist $\varphi'(s)=0$, so bedeutet dieses, dass die Gleichung $\varphi=0$ zwei gleiche Wurzeln $s$ hat. Der zugehörige Werth von $z$ bestimmt eine Ebene des ersten Büschels (17.), welcher die Curve in $k$ Punkten schneidet; von diesen müssen zwei zusammenfallen, damit von den zugehörigen Ebenen des Büschels $s$ zwei identisch werden. Es können nun hiebei folgende Fälle eintreten. Entweder ist die betreffende Ebene des Büschels $z$ Tangentenebene der Curve; oder die Ebene $s$ schneidet sich mit der Ebene $z$ in einer Geraden, welche die Curve in zwei verschiedenen Punkten trifft; oder endlich die Curve hat einen Doppelpunkt, was eine Berührung der gegebenen Oberflächen oder einen Knotenpunkt in einer von beiden anzeigt. In den beiden letzten Fällen gehört ebensowohl ein Doppelwerth von $s$ zu einem Werthe von $z$, als umgekehrt ein Doppelwerth von $z$ zu einem Werthe von $s$; in beiden Fällen müssen also sowohl $\varphi'(s)$ als $\varphi'(z)$ verschwinden. Schliessen wir nur diejenigen Fälle aus, in denen der Doppelpunkt in einen Rückkehrpunkt übergeht und also auch noch $\varphi''(ss)\,\varphi''(zz)-\varphi''(sz)^{2}=0$.

Die Zahl $\mathrm{w}$ bedeutet also offenbar den \emph{Rang} der Curve, d.\ h.\ die Zahl derjenigen Tangenten der Curve, welche eine gegebene Gerade treffen. So oft nämlich die Axe des Büschels $s$ von einer Tangente der Curve getroffen wird, ist eine Ebene des Büschels Tangentenebene der Curve, und also für dieselbe $\varphi'(s)=0$, ohne dass $\varphi'(z)$ verschwindet. Nach Herrn \emph{Salmon} (Geometry of three dimensions p.\ 237) bestimmt sich diese Zahl aus der Formel
\[
\mathrm{w} = k(k-1)-2d,
\]
wo $d$ die Anzahl der von einem Punkte ausgehenden Strahlen bedeutet, welche der Curve zweimal begegnen, mag diese zweifache Begegnung nun, wie bei einem Doppelpunkte, in demselben, oder mag sie in verschiedenen Punkten stattfinden.

So hat man also endlich
\[
p = \frac{(k-1)(k-2)}{2}-d
\]
für die charakteristische Zahl der der Curve zugehörigen \emph{Abel}schen Functionen.

Ich werde im Folgenden voraussetzen, dass die betrachtete Curve der vollständige und nicht in Theile zerlegbare Durchschnitt zweier Oberflächen

\origpage{221}

sei, welche sich nicht berühren; dass die Curve keinerlei Rückkehrpunkte und Doppelpunkte besitze ${}^{*)}$. In diesem Falle wird die Ordnung der Curve $=mn$, und nach Herrn \emph{Salmon} (l.\ c.\ p.\ 249) der Rang derselben
\[
\mathrm{w} = mn\,(m+n-2),
\]
also
\[
p = \frac{mn\,(m+n-4)}{2}+1;
\]
eine Zahl welche, übereinstimmend mit dem Früheren, für $m=1$ in $\dfrac{n-1.n-2}{2}$ übergeht.

Es ist leicht, die $p$ überall endlichen Integrale zu bilden, welche diesem Falle entsprechen. Aus den Gleichungen
\[
\begin{aligned}
u_1 x_1+u_2 x_2+u_3 x_3+u_4 x_4 &= 0,\\
v_1 x_1+v_2 x_2+v_3 x_3+v_4 x_4 &= 0,\\
u_1\,dx_1+u_2\,dx_2+u_3\,dx_3+u_4\,dx_4 &= 0,\\
v_1\,dx_1+v_2\,dx_2+v_3\,dx_3+v_4\,dx_4 &= 0
\end{aligned}
\qquad
\left(u_i = \frac{\partial u}{\partial x_i}, \quad v_i = \frac{\partial v}{\partial x_i}\right)
\]
ergeben sich folgende, in denen $d\mu$ einen unbestimmten Factor bedeutet:
\[
\begin{aligned}
x_1\,dx_2-x_2\,dx_1 &= (u_3 v_4-u_4 v_3)\,d\mu, & x_2\,dx_3-x_3\,dx_2 &= (u_1 v_4-u_4 v_1)\,d\mu,\\
x_1\,dx_3-x_3\,dx_1 &= (u_4 v_2-u_2 v_4)\,d\mu, & x_3\,dx_4-x_4\,dx_3 &= (u_1 v_2-u_2 v_1)\,d\mu,\\
x_1\,dx_4-x_4\,dx_1 &= (u_2 v_3-u_3 v_2)\,d\mu, & x_4\,dx_2-x_2\,dx_4 &= (u_1 v_3-u_3 v_1)\,d\mu.
\end{aligned}
\]
Bezeichnen also $a_1$, $a_2$, $\ldots$, $b_1$, $b_2$, $\ldots$ ganz beliebige Grössen, so ist der Ausdruck
\[
\frac{\Sigma \pm a_1 b_2 x_3\,dx_4}{\Sigma a_i u_i . \Sigma b_i v_i - \Sigma a_i v_i . \Sigma b_i u_i} = d\mu,
\]
demnach von den Werthen dieser Grössen völlig unabhängig. Derselbe ist von der $-(m+n-4)^{\text{ten}}$ Dimension; das Integral
\[
\Psi = \int \frac{\Theta . \Sigma \pm a_1 b_2 x_3\,dx_4}{\Sigma a_i u_i . \Sigma b_i v_i - \Sigma a_i v_i . \Sigma b_i u_i}
\]
lässt sich also mit Hülfe von $u=0$, $v=0$ als das Integral eines Differentialausdrucks mit einer einzigen Variablen darstellen, sobald $\Theta$ eine homogene Function der Ordnung $(m+n-4)$ ist. Unter den gemachten Voraussetzungen bleibt dieses

\textbf{*)} Eine solche Curve soll bisweilen der Kürze wegen eine Curve $mn^{\text{ter}}$ Ordnung genannt werden. Die Natur der Curve hängt dann nicht blos von der Ordnungszahl, sondern auch von ihrer Zerlegung ab, d.\ h.\ von den Ordnungen der Flächen, welche sich in ihr durchschneiden.

\origpage{222}

Integral auch überall endlich, und es ist nur noch nöthig sich zu überzeugen, ob wirklich $p$ Integrale durch verschiedene Annahme von $\Theta$ erhalten werden. Bemerken wir hiezu, dass $\Theta$ mit Hülfe der Gleichungen $u=0$, $v=0$ umgeformt werden kann, ohne dass das Integral sich ändert. Man kann also, wenn $U$, $V$ zwei Functionen $(m-4)^{\text{ter}}$ und $(n-4)^{\text{ter}}$ Ordnung bedeuten, statt $\Theta$ die Function
\[
\Theta+Uu+Vv
\]
setzen, und die Constanten von $U$, $V$ so bestimmen, dass ebensoviel Constante in $\Theta$ verschwinden. Demnach ist die Zahl der in $\Theta$ enthaltenen willkürlichen Constanten noch
\[
\begin{aligned}
&\tfrac{1}{6}\{(m+n-1)(m+n-2)(m+n-3)-(m-1)(m-2)(m-3)-(n-1)(n-2)(n-3)\}\\
&\qquad = \frac{mn\,(m+n-4)}{2}+1 = p.
\end{aligned}
\]
Man erhält also wirklich weder mehr noch weniger als $p$ unabhängige, überall endliche Integrale, indem man diesen Constanten alle möglichen Werthe beilegt.

\section*{§.~12.}

\subsection*{Der Abelsche Satz auf den Schnitt algebraischer Flächen angewandt.}

Es sei nun $w=0$ eine Oberfläche $k^{\text{ter}}$ Ordnung. Betrachten wir den Ausdruck $\Sigma\Psi$ ausgedehnt über alle Schnittpunkte der Fläche $w=0$ mit der gegebenen Curve. Führen wir bei der Integration wieder die Coefficienten einer Function $w$ durch die Gleichung $w=0$ als Variable ein, so dass erst in den oberen Grenzen $w$ in die gegebene Fläche übergehe. Bezeichnen wir durch $\delta w$, was aus $w$ durch Variation der Coefficienten entsteht, so ist
\[
\Sigma w_i\,dx_i = -\delta w.
\]

Multipliciren wir also unter dem Integralzeichen von $\Psi$ mit der Determinante
\[
\Sigma \pm c_1 u_2 v_3 w_4,
\]
so erhalten wir im Zähler:
\[
\Sigma \pm a_1 b_2 x_3\,dx_4 . \Sigma \pm c_1 u_2 v_3 w_4 =
\begin{vmatrix}
\Sigma a_i c_i & \Sigma b_i c_i & \Sigma c_i x_i & \Sigma c_i\,dx_i\\
\Sigma a_i u_i & \Sigma b_i u_i & 0 & 0\\
\Sigma a_i v_i & \Sigma b_i v_i & 0 & 0\\
\Sigma a_i w_i & \Sigma b_i w_i & 0 & -\delta w
\end{vmatrix}
\]
\[
= -\delta w . \Sigma c_i x_i . \{\Sigma a_i u_i \Sigma b_i v_i - \Sigma b_i u_i \Sigma a_i v_i\},
\]

\end{document}
